\documentclass[11pt]{article}
\usepackage[margin=1in]{geometry}
\usepackage{amsmath,amssymb,amsthm,mathtools}
\usepackage{enumitem}
\usepackage[hidelinks]{hyperref}
\sloppy\emergencystretch=4em

\newtheorem{theorem}{Theorem}[section]
\newtheorem{lemma}[theorem]{Lemma}
\newtheorem{proposition}[theorem]{Proposition}
\newtheorem{corollary}[theorem]{Corollary}
\theoremstyle{definition}
\newtheorem{definition}[theorem]{Definition}
\newtheorem{inpt}{Input}
\theoremstyle{remark}
\newtheorem{remark}[theorem]{Remark}
\newtheorem*{remarkun}{Remark}% unnumbered, so that the numbering of the other results is unaffected

\newcommand{\R}{\mathbb R}
\newcommand{\C}{\mathbb C}
\newcommand{\Z}{\mathbb Z}
\newcommand{\N}{\mathbb N}
\newcommand{\dd}{\,\mathrm d}
\newcommand{\dep}{\delta}
\newcommand{\sq}{\mathrm{sq}}
\newcommand{\avg}{\fint}
\newcommand{\Wneu}{W^{\mathrm{neu}}}
\newcommand{\Wtri}{W_{\triangle}}
\newcommand{\loc}[1]{\texttt{#1}}
\newcommand{\lc}{\ell_{\rm c}}
\newcommand{\lt}{\ell_{\rm t}}
\DeclareMathOperator{\dist}{dist}
\DeclareMathOperator{\supp}{supp}
\DeclareMathOperator{\Div}{div}
\def\fint{\mathop{\,\rlap{--}\!\!\int}\nolimits}

\title{Screening by explicit fields, the projection onto a square torus,\\
and the lower bound for the jellium energy}
\author{Scott Armstrong \and Sylvia Serfaty}
\date{}

\begin{document}
\maketitle

\begin{abstract}
We prove, with explicit constants, that the jellium energy $\Wneu(E)$ of Chapter~11 of Serfaty's
\emph{Lectures on Coulomb and Riesz gases}~\cite{Lectures} is at least $W_\triangle/(2\pi)$ for every field
$E$ in $\mathrm{elec}(C)$, where $W_\triangle$ is the constant of the preprint~\cite{Preprint}, and that the
energy of the configuration $\Lambda_\triangle$ is the smallest one.  The route avoids Neumann problems and is
built only from explicit fields.  The ingredients are: the Lipschitz pairing and Lemma~4.13 of the Lectures in field
form (\S2); a density lemma and the identification of the energy of Definition~11.4 with the neutral energy, with
$C_0^\ast=\log2+\pi/32$ (\S3); a screening theorem for fields on a square, built only from
explicit fields (a rectangle solver obtained by one-dimensional integration, radial point-to-disc fields, a
Lipschitz cut, pinwheel columns cut at integer masses, flux cells and an equal-mass filler), with an exact weak
pairing for the interaction between kept singularities and new fields (\S4); the projection of a zero-flux field
on the square of side $\sqrt N$ onto the torus energy of $\C/\sqrt N\Z[i]$ (\S5); and the assembly (\S6).
\S7 derives the screening input of the sphere note from the screening theorem.  Every constant was also checked
numerically.  The only inputs are Theorem~1.2 of the preprint, the triangular field, results already in the
formalization, Mathlib, and one classical fact (removable singularities of bounded harmonic
functions, Input~\ref{in:removable}).
\end{abstract}

\section*{Introduction}
Theorem~1.2 of the preprint~\cite{Preprint} is a lower bound for the energy of $n$ points on the square torus
$\C/\sqrt n\Z[i]$.  To use it for the planar problem one needs the jellium energy of Chapter~11 of the
Lectures: for a locally finite configuration $C$ and a field $E\in\mathrm{elec}(C)$, the cube energy
$F^{\square_\ell}(E)$ is the energy of $E$ in the square of side $\ell$ with the self-energy of each charge
removed (Definition~\ref{def:objects}), and the neutral energy $\Wneu(E)$ is $\limsup_\ell\ell^{-2}F^{\square_\ell}(E)$ when
the density of charges is $1$ and $+\infty$ otherwise.  The lower bound for $\Wneu$ is deduced in the Lectures from
Lemma~11.2, Proposition~11.3, Lemma~11.9, Proposition~11.13 and the screening of Chapter~7 and Appendix~A, which
solves Neumann problems.  We give a different proof, which in the formalization is the proof of the following
statements:
\begin{enumerate}[label=(\roman*)]
\item \emph{the jellium lower bound} (Theorem~\ref{thm:R2}): $W_\triangle/2\pi\le\Wneu(E)$ for every $C$ and every
$E\in\mathrm{elec}(C)$, in the extended reals
(\texttt{SandierSerfaty.div\_two\_pi\_triangularEnergy\_le\_neutralEnergy});
\item \emph{the bridge to Definition~11.4} (Theorem~\ref{thm:R2prime}): the energy $W_{C_0}$ of Definition~11.4 of the
Lectures equals $\Wneu$ for every $C_0\ge\log2+\pi/32$
(\texttt{SandierSerfaty.exists\_jelliumEnergy\_eq\_neutralEnergy});
\item \emph{the configuration lower bound} (Theorem~\ref{thm:R3}): the energy of the triangular lattice is
$W_\triangle/2\pi$ and is the smallest among all configurations
(\texttt{SandierSerfaty.configEnergy\_triangularLattice}).
\end{enumerate}
Statement~(i) is used by the comparison note, which deduces Theorem~1.1 of the preprint from it.

The proof of (i) screens a given field on a square: after discarding the charges near the boundary, the
field is modified in a boundary layer so that it has zero normal flux and exactly the right number of charges
$N=\ell^2$ in the square of side $\ell=\sqrt N$ (Theorem~\ref{thm:I4}), at a cost $o(\ell^2)$ in energy when the energy
density is bounded.  The modification uses only explicit fields.  The screened field is then periodized, and its
energy is compared with the torus energy of the corresponding $N$ points (Theorem~\ref{thm:I5}), to which
Theorem~1.2 of the preprint applies.  Statement~(ii) needs only the density lemma of \S3.  \S\ref{sec:sphere} shows
that the screening theorem also gives the screening input (Input~6.3) of the sphere note.  The formalization
is in \texttt{SandierSerfaty/Screening/}, and the Lean name of each main statement is given in its title.

\tableofcontents

%%%%%%%%%%%%%%%%%%%%%%%%%%%%%%%%%%%%%%%%%%%%%%%%%%%%%%%%%%%%%%%%%%%%%%%%%%%%%%%%%%%%%%%%%%%%%%%%%%%%%%%%%
\section{Setting, normalizations and inputs}\label{sec:setting}

\subsection{Conventions}
We work in $\C$ with the Euclidean norm $|x|$; $|x|_\infty=\max(|\operatorname{Re}x|,|\operatorname{Im}x|)$.
$B(x,r)$ is the open disc, $\bar B(x,r)$ the closed disc, $\sigma_{x,r}$ the normalized arclength measure on
$\partial B(x,r)$ and $\avg_{\partial B(x,r)}\varphi=\int\varphi\dd\sigma_{x,r}$ (Lean: \texttt{Real.circleAverage}).
For $c\in\C$, $s\ge0$ the closed square of side $s$ is
\[
\sq(c,s)=\{x:|x-c|_\infty\le s/2\}\qquad(\text{Lean: } \texttt{Set.Icc}\times_{\C}\texttt{Set.Icc}),
\]
and $\texttt{cube}\,\ell=\square_\ell=\sq(0,\ell)$. For a fixed square $Q=\sq(c,t)$ the
\emph{(signed) depth} is $\dep(x)=t/2-|x-c|_\infty$. It is $1$-Lipschitz for $|\cdot|_\infty$, hence for $|\cdot|$, and
\begin{equation}\label{eq:depth-sets}
\sq(c,t-2s)=\{\dep\ge s\},\qquad \operatorname{int}\sq(c,t-2s)=\{\dep>s\},\qquad
\dist(x,\partial Q)=\dep(x)\ (x\in Q).
\end{equation}
(The last identity: for $x\in Q$ the nearest boundary point is the orthogonal projection on the nearest side.)
A function $\varphi$ is a \emph{Lipschitz test function on $U$} if it is Lipschitz and has compact support in the
open set $U$; its gradient is the a.e.\ gradient (Rademacher), which is what Mathlib's \texttt{gradient} returns
a.e.\ (it is $0$ at points of non-differentiability, a null set).

\subsection{\texorpdfstring{The objects of the statements ($d=2$, $s=0$, $c_{2,0}=2\pi$, $g=-\log|\cdot|$)}{The objects of the statements}}
All definitions below are the Lean definitions that enter the statements (in the namespace \texttt{SandierSerfaty}); we
restate them in mathematical notation.

\begin{definition}[definitions of the statements]\label{def:objects}
\begin{enumerate}[label=(\alph*)]
\item \texttt{truncLog}: $f_\eta(x)=\max(\log\eta-\log|x|,0)$ ($\eta>0$; Lectures (4.1.18));
  $\nabla f_\eta(x)=-x|x|^{-2}\mathbf 1_{B(0,\eta)}(x)$ for $|x|\ne\eta$, $x\ne0$.
\item \texttt{elec}: $E\in\mathrm{elec}(C)$ iff $C$ is locally finite, $E\in L^1_{\rm loc}(\C;\C)$ and
  $\int E\cdot\nabla\varphi=2\pi\bigl(\sum_{p\in C}\varphi(p)-\int\varphi\bigr)$ for all real
  $\varphi\in C^\infty_c(\C)$ (Lectures, Def.~11.1); i.e.\ $-\Div E=2\pi(\sum_C\delta_p-1)$.
\item \texttt{truncField}: $E_\eta=E-\sum_{p\in C}\nabla f_{\eta_p}(\cdot-p)$ for a radius family $\eta$
  (Lectures (11.2.4)).
\item \texttt{nnRadius}: $r_p=\frac14\inf\bigl(\{1\}\cup\{|q-p|:q\in C\setminus\{p\}\}\bigr)$; \texttt{cubeRadius}:
  $\tilde r_p=r_p$ if $p\in\square_\ell$, $\dist(p,\partial\square_\ell)\ge2$; $\tilde r_p=\frac14$ if
  $p\notin\square_\ell$ or $\dist(p,\partial\square_\ell)\le1$; linear interpolation in between
  (Lectures (11.2.2)--(11.2.3)). Always $0<\tilde r_p\le\frac14$ for locally finite $C$.
\item \texttt{cubeEnergy}: $F^{\square_\ell}(E)=\frac1{4\pi}\int_{\square_\ell}|E_{\tilde r}|^2
  -\frac12\sum_{p\in C\cap\square_\ell}(-\log\tilde r_p)-\sum_{p\in C\cap\square_\ell}\int f_{\tilde r_p}(x-p)\dd x$,
  the field integral a lower Lebesgue integral (Lectures (11.2.6)).
\item \texttt{neutralEnergy}: $\Wneu(E)=\limsup_{\ell\to\infty}\ell^{-2}F^{\square_\ell}(E)$ (real $\ell$) if
  $n(\ell)/\ell^2\to1$, where $n(\ell)=\#(C\cap\square_\ell)$, and $+\infty$ otherwise.
  \texttt{jelliumEnergy}: $W_{C_0}(E)=\limsup_R R^{-2}[F^{\square_R}(E)+C_0(n(R)-R^2)]$ (Def.~11.4 of the Lectures). \texttt{configEnergy}: $\mathbb W(C)=\inf\{\Wneu(E):E\in\mathrm{elec}(C),\ E=\nabla H
  \text{ distributionally},\ H\in L^1_{\rm loc}\}$.
\item Torus: $\Gamma_N=\texttt{squareLattice}\,N=\sqrt N\,\Z[i]$; for $a:\mathrm{Fin}\,N\to\C$ with distinct classes,
  $H=\texttt{torusPotential}\,\Gamma\,a$ is the mean-zero $\Gamma$-periodic solution of $-\Delta H=2\pi(\sum_i\delta_{a_i+\Gamma}-1)$
  in the classical sense of \texttt{IsTorusPotential} (smooth off the charges, $\Delta H=2\pi$ there,
  $H+\log|\cdot-a_i|$ has a limit at $a_i$, $\int_TH=0$, $H=0$ at the charges), and
  \[
  \texttt{torusEnergy}\,\Gamma\,a=\lim_{\eta\downarrow0}\Bigl[\tfrac12\int_{T\setminus\bigcup_{i,\gamma}B(a_i+\gamma,\eta)}|\nabla H|^2+\pi N\log\eta\Bigr],
  \]
  $T=\texttt{torusDomain}\,\Gamma$ (an a.e.\ fundamental domain). $W_\triangle=\texttt{triangularEnergy}$ is
  this energy for the one-point triangular torus, and $E_\triangle=\texttt{triangularField}=\nabla H_\triangle$.
\end{enumerate}
\end{definition}

\begin{remark}[sign]\label{rem:sign}
The sign of the Lectures is used throughout: $-\Div E=2\pi(\sum\delta_p-m)$, so near a charge $E\approx\nabla(-\log|x-p|)=-(x-p)/|x-p|^2$,
and $E_\eta=E-\sum\nabla f_{\eta_p}(\cdot-p)$ removes this singularity. The comparison note uses the sign of the paper,
$\tilde E=-E$; all its statements transfer by $E\mapsto-E$.
\end{remark}

\begin{definition}[charge fields; local energy]\label{def:local}
Let $U\subset\C$ be open, $m\in\R$, $C\subset\C$ locally finite.
\begin{enumerate}[label=(\alph*)]
\item $E$ is a \emph{charge field on $U$ with density $m$} (\texttt{IsChargeField U m C E}) if $E\in L^1_{\rm loc}(U)$ and
  $\int E\cdot\nabla\varphi=2\pi(\sum_{p\in C}\varphi(p)-m\int\varphi)$ for $\varphi\in C_c^\infty(U)$.
\item For a closed square $Q$ and a finite $Y\subset\C$, $E$ is a \emph{Neumann charge field on $Q$}
  (\texttt{IsNeumannChargeField Q m Y E}) if $E\in L^1(Q)$ and
  $\int_QE\cdot\nabla\varphi=2\pi(\sum_{y\in Y}\varphi(y)-m\int_Q\varphi)$ for every real $\varphi\in C^\infty_c(\C)$
  (restrictions of global test functions: this encodes $-\Div E=2\pi(\sum\delta_y-m)$ in $Q^\circ$ and zero normal flux).
\item $\tau_m(r)=\pi\log r-\pi^2mr^2$, $c_m=-\tau_m(\tfrac14)=\pi\log4+\pi^2m/16$
  ($c_1=4.972022\ldots$).
\item For a radius family $\alpha>0$ and a bounded Borel set $D$:
  \[
  G_D(\alpha)[E]=\texttt{locEnergy}\,m\,C\,\alpha\,E\,D=\tfrac12\int_D|E_\alpha|^2+\sum_{p\in C\cap D}\tau_m(\alpha_p)\in(-\infty,+\infty],
  \]
  the integral being a lower Lebesgue integral (so $G_D=+\infty$ iff $E_\alpha\notin L^2(D)$).
\end{enumerate}
\end{definition}

\begin{lemma}[the $2\pi$ normalization]\label{lem:norm}
(a) If $E\in\mathrm{elec}(C)$ then $E$ is a charge field on every open $U$ with density $1$.
(b) For locally finite $C$, every $E$ and $\ell>0$: $2\pi F^{\square_\ell}(E)=G_{\square_\ell}(\tilde r)[E]$ in $(-\infty,+\infty]$,
with $m=1$ (\texttt{two\_pi\_mul\_cubeEnergy}).
\end{lemma}
\begin{proof}
(a) Restriction of test functions. (b) $\int f_\eta=\pi\eta^2/2$ (\texttt{PolarMean.integral\_truncLog}). Multiplying the definition by $2\pi$ gives
$\frac12\int|E_{\tilde r}|^2+\sum[\pi\log\tilde r_p-2\pi\cdot\pi\tilde r_p^2/2]=\frac12\int|E_{\tilde r}|^2+\sum\tau_1(\tilde r_p)$;
in $(-\infty,+\infty]$ multiplication by $2\pi>0$ distributes over $(+\infty)+(\text{real})$.
\end{proof}

\subsection{Inputs}\label{sec:inputs}
\begin{inpt}[Theorem~1.2 of the preprint: the lower bound on square tori]\label{in:R1}
\texttt{mul\_triangularEnergy\_le\_torusEnergy}:
for $n\ge2$ and $a:\mathrm{Fin}\,n\to\C$ with pairwise distinct classes mod $\Gamma_n$,
$nW_\triangle\le\texttt{torusEnergy}(\Gamma_n,a)$.  (Proposition~\ref{prop:I6} below takes this
statement as its hypothesis.)
\end{inpt}
\begin{inpt}[the triangular field]\label{in:Etri}
\texttt{triangularField\_spec}: $E_\triangle\in\mathrm{elec}(\Lambda_\triangle)$ and $\Wneu(E_\triangle)=W_\triangle/2\pi$.
\end{inpt}
\begin{inpt}[facts about the torus potential]\label{in:main}
\texttt{gradient\_torusPotential\_mem\_elec}: if $\operatorname{covol}\Gamma=N$ and the
classes of $a$ are distinct, $\nabla H\in\mathrm{elec}(a+\Gamma)$; \texttt{isTorusPotential\_torusPotential};
\texttt{PolarMean.integral\_truncLog}; \texttt{PolarDivergence.integral\_exterior\_div}; \texttt{torusDomain\_spec},
\texttt{PeriodicIntegral.integral\_torusDomain\_add}; \texttt{Torus.exists\_local\_form\_ball} (one charge).
\end{inpt}
\begin{inpt}[classical; removable singularity]\label{in:removable}
A harmonic function on a punctured disc $B(a,s)\setminus\{a\}$ that is bounded near $a$ extends to a harmonic function on
$B(a,s)$. (Proof: write $h=\operatorname{Re}F+\kappa\log|z-a|$ with $F$ holomorphic on the punctured disc and
$\kappa=\frac1{2\pi}\int_{\partial B(a,r)}\partial_\nu h$; boundedness forces $\kappa=0$; then $e^{F}$ is bounded and bounded away from $0$
near $a$, so by Riemann's theorem $e^F$ and hence $F=\log e^F$ (locally) extend holomorphically. In the formalization this is
used in \texttt{Torus.exists\_local\_form\_ball}.)
\end{inpt}
\textbf{Mathlib:} Rademacher (\texttt{LipschitzWith.ae\_differentiableAt},
\texttt{LipschitzWith.integral\_lineDeriv\_mul\_eq}), mollifiers (\texttt{ContDiffBump},
\texttt{ContDiffBump.ae\_convolution\_tendsto\_right\_of\_locallyIntegrable}, \texttt{HasCompactSupport.contDiff\_convolution\_left}),
Fubini/Tonelli (\texttt{integral\_prod}, \texttt{lintegral\_prod}), \texttt{intermediate\_value\_Icc},
\texttt{IsAddFundamentalDomain}, EReal $\limsup$ API (\texttt{le\_limsup\_of\_frequently\_le}).

%%%%%%%%%%%%%%%%%%%%%%%%%%%%%%%%%%%%%%%%%%%%%%%%%%%%%%%%%%%%%%%%%%%%%%%%%%%%%%%%%%%%%%%%%%%%%%%%%%%%%%%%%
\section{Truncation calculus, the Lipschitz pairing and Lemma 4.13}\label{sec:mono}

\begin{lemma}[the functions $f_\eta$; Lemma~2.2 of the comparison note, in the sign of the Lectures]\label{lem:f}
Let $0<\alpha\le\beta$, $f_{\alpha,\beta}=f_\alpha-f_\beta$.
\begin{enumerate}[label=(\alph*)]
\item $f_\eta\in W^{1,1}_{\rm loc}$, $\nabla f_\eta=-x|x|^{-2}\mathbf 1_{B(0,\eta)}$ a.e.; $\int_\C f_\eta=\pi\eta^2/2$.
\item $f_{\alpha,\beta}$ is $\frac1\alpha$-Lipschitz, $=\log(\alpha/\beta)$ on $\bar B(0,\alpha)$, $=\log(|x|/\beta)$ for
  $\alpha\le|x|\le\beta$, $=0$ for $|x|\ge\beta$; so $\log(\alpha/\beta)\le f_{\alpha,\beta}\le0$ and
  $\int|f_{\alpha,\beta}|=\frac\pi2(\beta^2-\alpha^2)$; in particular $\int|f_{\alpha,1/4}|\le\pi/32$.
\item For every Lipschitz $\varphi$ with compact support: $\int\nabla f_\eta\cdot\nabla\varphi=2\pi\bigl(\varphi(0)-\avg_{\partial B(0,\eta)}\varphi\bigr)$,
  i.e.\ $-\Delta f_\eta=2\pi(\delta_0-\sigma_{0,\eta})$; hence $\int\nabla f_{\alpha,\beta}\cdot\nabla\varphi=2\pi\int\varphi\dd(\sigma_{0,\beta}-\sigma_{0,\alpha})$.
\item $\int f_{\alpha,\beta}\dd(\sigma_{0,\alpha}+\sigma_{0,\beta})=\log(\alpha/\beta)$.
\end{enumerate}
\end{lemma}
\begin{proof}
(a),(b),(d): the formulas and polar coordinates ($\int_0^\eta\log(\eta/r)2\pi r\dd r=\pi\eta^2/2$).
(c) For $\varphi\in C^\infty_c$ and $0<\epsilon<\eta$, the divergence theorem for $\varphi\,x|x|^{-2}$ on $B_\eta\setminus\bar B_\epsilon$
($\Div(x|x|^{-2})=0$ there; \texttt{PolarDivergence.integral\_exterior\_div}) gives
$\int_{B_\eta\setminus B_\epsilon}x|x|^{-2}\cdot\nabla\varphi=2\pi\avg_{\partial B_\eta}\varphi-2\pi\avg_{\partial B_\epsilon}\varphi$;
let $\epsilon\to0$ and multiply by $-1$. For Lipschitz $\varphi$ apply Lemma~\ref{lem:moll} with $Y=\nabla f_\eta\in L^1_{\rm loc}$ and the measure
$2\pi(\delta_0-\sigma_{0,\eta})$.
\end{proof}

\begin{remarkun}[continuous representative of $f_{\alpha,\beta}$ as a test function]\label{rem:contrep}
Whenever $f_{\alpha,\beta}(\cdot-p)$ is used below as a Lipschitz test function (Lemma~\ref{lem:f}(b),(c); Proposition~\ref{prop:mono}; Lemma~\ref{lem:sign}), it
denotes its \emph{continuous representative}
\[
f_{\alpha,\beta}(x)=\max\bigl(\log(\alpha/\beta),\min(0,\log(|x|/\beta))\bigr)=\log\frac{\max(\alpha,\min(|x|,\beta))}{\beta}\qquad(x\ne0),
\]
extended by $\log(\alpha/\beta)$ at $x=0$ (the second form is the one that needs no value of $\log$ at $0$). This is exactly the piecewise formula of Lemma~\ref{lem:f}(b).
It differs from the difference $f_\alpha-f_\beta$ of the \texttt{truncLog} functions at the single point $0$ only, because the Lean value
$\texttt{truncLog}\,\eta\,0=\max(\log\eta,0)$ ($\log0=0$) is not the limit at the centre. Hence a.e.\ gradients, integrals and circle averages
are unchanged (Lean: \texttt{circleAverage\_truncLog\_sub\_eq}), while the Lipschitz and compact-support hypotheses of Lemma~\ref{lem:moll} hold for the representative.
\end{remarkun}

\begin{lemma}[extension to Lipschitz test functions]\label{lem:moll}
Let $U$ be open, $Y\in L^1_{\rm loc}(U;\C)$ and $\nu$ a real Radon measure on $U$ (finite total variation on compacts)
with $\int Y\cdot\nabla\varphi=\int\varphi\dd\nu$ for all $\varphi\in C^\infty_c(U)$. Then the same identity holds for every
Lipschitz test function $\varphi$ on $U$.
\end{lemma}
\begin{proof}
Let $K=\supp\varphi$, $3\epsilon_0=\dist(K,\C\setminus U)$ ($=\infty$ if $U=\C$), $K'=K+\bar B(0,\epsilon_0)$, and
$\varphi_\epsilon=\varrho_\epsilon*\varphi$ for a standard mollifier, $\epsilon<\epsilon_0$. Then $\varphi_\epsilon\in C_c^\infty(U)$,
$\supp\varphi_\epsilon\subset K'$, $\varphi_\epsilon\to\varphi$ uniformly, $\nabla\varphi_\epsilon=\varrho_\epsilon*\nabla\varphi$ (the weak gradient
of a Lipschitz function is its a.e.\ gradient: \texttt{integral\_lineDeriv\_mul\_eq}), so $|\nabla\varphi_\epsilon|\le\operatorname{Lip}\varphi$ and
$\nabla\varphi_\epsilon\to\nabla\varphi$ a.e.\ (Lebesgue differentiation for the convolution). Dominated convergence with
dominating function $\operatorname{Lip}(\varphi)|Y|\mathbf 1_{K'}\in L^1$ on the left; uniform convergence on $K'$ and
$|\nu|(K')<\infty$ on the right.
\end{proof}

\begin{proposition}[Lipschitz pairing; Lectures (4.1.22), Lemma~2.4 of the comparison note]\label{prop:pairing}
Let $E$ be a charge field on $U$ with density $m$ (charges $C$), and $\alpha$ a positive radius family with
$\sup_C\alpha<\infty$. For every Lipschitz test function $\varphi$ on $U$:
\begin{align}
\int E\cdot\nabla\varphi&=2\pi\Bigl(\sum_{p\in C}\varphi(p)-m\int\varphi\Bigr),\label{eq:pair0}\\
\int E_\alpha\cdot\nabla\varphi&=2\pi\Bigl(\sum_{p\in C}\avg_{\partial B(p,\alpha_p)}\varphi-m\int\varphi\Bigr).\label{eq:pair1}
\end{align}
All sums are finite.
\end{proposition}
\begin{proof}
\eqref{eq:pair0}: Lemma~\ref{lem:moll} with $\nu=2\pi(\sum_{C\cap U}\delta_p-m\,\mathcal L^2)$ (Radon on $U$ since $C$ is locally finite).
\eqref{eq:pair1}: subtract Lemma~\ref{lem:f}(c) at the points $p$ with $\bar B(p,\alpha_p)\cap\supp\varphi\ne\emptyset$ (finitely many,
as $\alpha$ is bounded and $C$ locally finite); for the other $p$ both $\varphi(p)$ and the circle average vanish.
The points $p\in C\setminus U$ are allowed: they contribute $0$ to \eqref{eq:pair0} and their circle averages to \eqref{eq:pair1}.
\end{proof}

\begin{proposition}[Lemma 4.13 of the Lectures in field form, exact; Prop.~3.1 of the comparison note]\label{prop:mono}
Let $E$ be a charge field on $U$ with density $m$ (charges $C$), $D$ a bounded Borel set, and $0<\alpha\le\beta$ radius families
on $C$ with $\sup\beta<\infty$. Let $I=\{i\in C:\alpha_i<\beta_i\}$ and assume $\bar B(i,\beta_i)\subset U\cap D$ for $i\in I$ and
$E_\alpha\in L^2(D)$. Then $I$ is finite, $E_\beta\in L^2(D)$, and
\begin{equation}\label{eq:mono}
G_D(\beta)-G_D(\alpha)=\mathcal I(\alpha,\beta):=\pi\sum_{i\in I}\ \sum_{q\in C\setminus\{i\}}\int f_{\alpha_i,\beta_i}(x-i)\dd\bigl(\sigma_{q,\alpha_q}+\sigma_{q,\beta_q}\bigr)(x)\ \le 0 .
\end{equation}
\end{proposition}
\begin{proof}
$I\subset D$ is bounded, so finite. Put $f_i=f_{\alpha_i,\beta_i}(\cdot-i)$: Lipschitz, $\le0$, support in $\bar B(i,\beta_i)\subset U\cap D$. By the definitions
$E_\beta=E_\alpha+\sum_{i\in I}\nabla f_i$ ($\nabla f_\alpha-\nabla f_\beta=\nabla f_{\alpha,\beta}$), so $E_\beta\in L^2(D)$ and
\[
T:=\int_D|E_\beta|^2-\int_D|E_\alpha|^2=2\sum_{i\in I}\int E_\alpha\cdot\nabla f_i+\sum_{i,j\in I}\int\nabla f_i\cdot\nabla f_j
\]
(the integrands vanish off $D$). By \eqref{eq:pair1}, $2\int E_\alpha\cdot\nabla f_i=4\pi\int f_i\dd\sigma_\alpha-4\pi m\int f_i$ with
$\sigma_\alpha=\sum_q\sigma_{q,\alpha_q}$. By Lemma~\ref{lem:f}(c) (translated),
$\int\nabla f_j\cdot\nabla f_i=2\pi\int f_i\dd(\sigma_{j,\beta_j}-\sigma_{j,\alpha_j})$. Since $\sigma_\beta=\sigma_\alpha+\sum_{j\in I}(\sigma_{j,\beta_j}-\sigma_{j,\alpha_j})$,
\[
T=2\pi\sum_{i\in I}\int f_i\dd(\sigma_\alpha+\sigma_\beta)-4\pi m\sum_{i\in I}\int f_i .
\]
The $q=i$ term of $\int f_i\dd(\sigma_\alpha+\sigma_\beta)$ is $\log(\alpha_i/\beta_i)$ (Lemma~\ref{lem:f}(d)); the others sum to
$o_i:=\sum_{q\ne i}\int f_i\dd(\sigma_{q,\alpha_q}+\sigma_{q,\beta_q})\le0$ (finitely many non-zero terms). By Lemma~\ref{lem:f}(a),
$-4\pi m\int f_i=2\pi^2m(\beta_i^2-\alpha_i^2)$. Hence $\frac12T=\sum_{i\in I}[\pi o_i-\pi\log(\beta_i/\alpha_i)+\pi^2m(\beta_i^2-\alpha_i^2)]$ while
$\sum_{C\cap D}[\tau_m(\beta_p)-\tau_m(\alpha_p)]=\sum_{i\in I}[\pi\log(\beta_i/\alpha_i)-\pi^2m(\beta_i^2-\alpha_i^2)]$.
Adding, $G_D(\beta)-G_D(\alpha)=\pi\sum_io_i$.
\end{proof}

\begin{corollary}[sign and equality; \texttt{truncInteraction\_nonpos}, \texttt{truncInteraction\_eq\_zero}]\label{cor:mono}
In Proposition~\ref{prop:mono}: (a) $G_D(\beta)\le G_D(\alpha)$. (b) If $\beta_i+\beta_q\le|i-q|$ for all $i\in I$, $q\in C\setminus\{i\}$, then
$G_D(\beta)=G_D(\alpha)$.
\end{corollary}
\begin{proof}
(a) $o_i\le0$. (b) The circles $\partial B(q,\alpha_q),\partial B(q,\beta_q)$ lie in $\bar B(q,\beta_q)$, which meets $\bar B(i,\beta_i)$ at most in
$\partial B(i,\beta_i)$, where $f_i=0$. So every $o_i=0$.
\end{proof}

\begin{remarkun}[hypotheses of the equality case]\label{rem:eqzero}
Corollary~\ref{cor:mono}(b) (\texttt{truncInteraction\allowbreak\_eq\allowbreak\_zero}) needs, besides $\beta_i+\beta_q\le|i-q|$, the hypotheses $0<\alpha_p$ and $\alpha_p\le\beta_p$ for all
$p\in C$ (those of \texttt{truncInteraction\allowbreak\_nonpos}); these are part of the standing assumption $0<\alpha\le\beta$ of Proposition~\ref{prop:mono},
and must be stated explicitly in the Lean statement. Without $\alpha\le\beta$ it is false: take $C=\{0,10\}$, $\alpha_0=1$, $\beta_0=2$, $\alpha_{10}=9.5$,
$\beta_{10}=1$; then $\beta_0+\beta_{10}\le10$ holds, but the circle of radius $9.5$ about $10$ passes within $0.5$ of $0$, where $f_{1,2}=-\log2<0$.
Every use below has $0<\alpha\le\beta$.
\end{remarkun}

\begin{corollary}[energy bound; \texttt{half\_lintegral\_quarter\_le}]\label{cor:EB}
Let $E\in\mathrm{elec}(C)$, $\ell>0$, $\Phi=E_{1/4}$ (constant family $\frac14$). Then, in $[0,+\infty]$,
\[
\tfrac12\int_{\square_\ell}|\Phi|^2\le2\pi F^{\square_\ell}(E)+c_1\,n(\ell),\qquad c_1=\pi\log4+\pi^2/16 .
\]
\end{corollary}
\begin{proof}
If $E_{\tilde r}\notin L^2(\square_\ell)$ the right side is $+\infty$. Otherwise apply Proposition~\ref{prop:mono} with $U=\C$, $D=\square_\ell$,
$\alpha=\tilde r$, $\beta\equiv\frac14$: $I=\{p:\tilde r_p<\frac14\}$ consists of points of $\square_\ell$ with $\dist(p,\partial\square_\ell)>1$, so
$\bar B(p,\frac14)\subset\square_\ell$. Then $G_{\square_\ell}(\tfrac14)\le G_{\square_\ell}(\tilde r)=2\pi F^{\square_\ell}$ (Lemma~\ref{lem:norm}) and
$G_{\square_\ell}(\frac14)=\frac12\int|\Phi|^2-c_1n(\ell)$.
\end{proof}

\begin{remark}[what is not needed]
The error in Lemma~11.2 of the Lectures (see the errata record of the repository) concerns the equality case for overlapping $\tilde r$-balls; it is not used: only
Corollary~\ref{cor:mono}(a) and the equality (b) for \emph{disjoint} balls appear below. Lemma~11.2, Prop.~11.3 and Lemma~11.9 of the Lectures are not used.
\end{remark}

\begin{lemma}[local count]\label{lem:localcount}
Let $E$ be a charge field on $U$ with density $m\ge0$ (charges $C$), $\Phi=E_{1/4}$, $x\in\C$, $r>0$, and $\bar B(x,r+\frac54)\subset U$ with
$\Phi\in L^2(B(x,r+\frac54))$. Then
\[
\#(C\cap\bar B(x,r))\le a(r,m)+b(r)\,\|\Phi\|_{L^2(B(x,r+5/4))},\qquad a(r,m)=m\pi(r+\tfrac54)^2,\quad b(r)=\sqrt{\tfrac{2r+3/2}{4\pi}} .
\]
\end{lemma}
\begin{proof}
$\varphi(y)=\min\{1,\max\{0,r+\frac54-|y-x|\}\}$ is $1$-Lipschitz, $=1$ on $\bar B(x,r+\frac14)$, supported in $\bar B(x,r+\frac54)\subset U$,
$\nabla\varphi$ supported in the annulus $r+\frac14\le|y-x|\le r+\frac54$ of area $\pi(2r+\frac32)$. All circle averages
$\avg_{\partial B(p,1/4)}\varphi$ are $\ge0$ and equal $1$ if $|p-x|\le r$. By \eqref{eq:pair1} and Cauchy--Schwarz,
$\#(C\cap\bar B(x,r))\le m\int\varphi+\frac1{2\pi}\int\Phi\cdot\nabla\varphi\le m\pi(r+\frac54)^2+\frac{\sqrt{\pi(2r+3/2)}}{2\pi}\|\Phi\|_{L^2(B(x,r+5/4))}$.
\end{proof}

%%%%%%%%%%%%%%%%%%%%%%%%%%%%%%%%%%%%%%%%%%%%%%%%%%%%%%%%%%%%%%%%%%%%%%%%%%%%%%%%%%%%%%%%%%%%%%%%%%%%%%%%%
\section{The density lemma and the bridge to Definition 11.4}\label{sec:density}

\begin{lemma}[grid counts]\label{lem:grid}
Let $\Z^2/2=\{x\in\C:2x\in\Z[i]\}$.
\begin{enumerate}[label=(\alph*)]
\item Every $y\in\C$ has $x\in\Z^2/2$ with $|y-x|_\infty\le\frac14$, hence $|y-x|\le r_1:=\sqrt2/4$.
\item A point of $\C$ lies in at most $M(R):=(\lfloor4R\rfloor+1)^2$ discs $B(x,R)$, $x\in\Z^2/2$. With $R_1=r_1+\frac54$,
  $R_2=r_2+\frac54$, $r_2=r_1+\frac12$: $M_1=M(R_1)=49$, $M_2=M(R_2)=81$.
\item A closed axis-parallel rectangle of sides $A,B$ contains at most $(2A+1)(2B+1)$ points of $\Z^2/2$.
\item For a square $Q=\sq(c,t)$ with depth $\dep$ and $1\le s_1<s_2$ with $t-2s_2\ge\frac12$:
  $\#\{x\in\Z^2/2:s_1<\dep(x)<s_2\}\le8t(2(s_2-s_1)+1)$.
\end{enumerate}
\end{lemma}
\begin{proof}
(a) Round each coordinate to $\frac12\Z$. (b),(c): the grid points in an open interval of length $2R$ (closed interval of length $A$)
number at most $\lfloor4R\rfloor+1$ ($\lfloor2A\rfloor+1$) per axis. (d) The set is contained in $\sq(c,L)\setminus\operatorname{int}\sq(c,L')$,
$L=t-2s_1$, $L'=t-2s_2$; a closed square of side $L$ contains at most $(2L+1)^2$ grid points and an open one of side $L'\ge\frac12$ at
least $(2L'-1)^2$; the difference is $4(L-L'+1)(L+L')\le8t(2(s_2-s_1)+1)$.
\end{proof}

\begin{proposition}[density; \texttt{tendsto\_ncard\_cube\_div\_sq}; compare Prop.~6.3 of the comparison note]\label{prop:density}
Let $E\in\mathrm{elec}(C)$, $\Phi=E_{1/4}$, and assume $\int_{\square_\ell}|\Phi|^2\le A\ell^2$ for all $\ell\ge\ell_1$ ($A\ge0$). Then for $\ell\ge\ell_1+7$,
\[
|n(\ell)-\ell^2|\le e(\ell)+8\ell,\qquad e(\ell):=\tfrac1{2\pi}\sqrt{(4\ell+4)A}\,(\ell+2)+40(\ell+1)a_1+b_1\sqrt{1960(\ell+1)A}\,(\ell+7),
\]
with $a_1=a(r_1,1)$, $b_1=b(r_1)$ (Lemma~\ref{lem:localcount}). In particular $e(\ell)=O(\ell^{3/2})$ and $n(\ell)/\ell^2\to1$.
\end{proposition}
\begin{proof}
Let $\varphi_\ell(y)=\min\{1,\max\{0,\ell/2+1-|y|_\infty\}\}$: $1$-Lipschitz, $\mathbf 1_{\square_\ell}\le\varphi_\ell\le\mathbf 1_{\square_{\ell+2}}$,
$\nabla\varphi_\ell$ supported in $\square_{\ell+2}\setminus\operatorname{int}\square_\ell$ (area $4\ell+4$).
By \eqref{eq:pair1} ($U=\C$, $m=1$) and Cauchy--Schwarz,
$|\sum_p\avg_{\partial B(p,1/4)}\varphi_\ell-\int\varphi_\ell|\le\frac1{2\pi}\sqrt{4\ell+4}\,\sqrt{A}(\ell+2)$.
If $|p|_\infty\le\ell/2-\frac14$ or $|p|_\infty\ge\ell/2+\frac54$ the circle average equals $\varphi_\ell(p)\in\{0,1\}$; the other points lie in
$\{\ell/2-\frac14<|p|_\infty<\ell/2+\frac54\}$, and each is within $r_1$ of a grid point $x$ with $\ell/2-\frac12<|x|_\infty<\ell/2+\frac32$;
by the count in the proof of Lemma~\ref{lem:grid}(d) (closed square of side $L=\ell+3$ minus open square of side $L'=\ell-1$)
there are at most $4(L-L'+1)(L+L')=40(\ell+1)$ such $x$, and $B(x,R_1)\subset\square_{\ell+7}$. Lemma~\ref{lem:localcount}, Cauchy--Schwarz and
Lemma~\ref{lem:grid}(b) bound their number by $40(\ell+1)a_1+b_1\sqrt{40(\ell+1)\cdot49\cdot A(\ell+7)^2}$. Each changes the sum by at most $1$. Hence
$|\sum_p\varphi_\ell(p)-\int\varphi_\ell|\le e(\ell)$. Now $n(\ell)\le\sum_p\varphi_\ell(p)\le n(\ell+2)$ and $\ell^2\le\int\varphi_\ell\le(\ell+2)^2$: so
$n(\ell)\le(\ell+2)^2+e(\ell)$ and $n(\ell)\ge\sum_p\varphi_{\ell-2}(p)\ge(\ell-2)^2-e(\ell-2)$.
\end{proof}

\begin{theorem}[bridge to Definition 11.4; \texttt{jelliumEnergy\_eq\_neutralEnergy}]\label{thm:R2prime}
Let $C_0^\ast=c_1/2\pi=\log2+\pi/32=0.79132\ldots$. For every $C_0\ge C_0^\ast$, every $C$ and every $E\in\mathrm{elec}(C)$:
$W_{C_0}(E)=\Wneu(E)$. Hence \texttt{exists\_jelliumEnergy\_eq\_neutralEnergy} holds with $C_0^\ast$ as the constant.
\end{theorem}
\begin{proof}
Write $a(R)=R^{-2}F^{\square_R}(E)\in(-\infty,+\infty]$ and $b(R)=C_0(n(R)/R^2-1)\in\R$ ($R>0$), so that $W_{C_0}=\limsup(a+b)$ (multiplication
by $R^{-2}>0$ distributes; no $\infty-\infty$ occurs since $F>-\infty$).

\emph{Case $n(R)/R^2\to1$.} Then $b(R)\to0$, and $\limsup(a+b)=\limsup a=\Wneu(E)$ in $[-\infty,+\infty]$ (a real sequence tending to $0$ does not change an
EReal $\limsup$).

\emph{Case $n(R)/R^2\not\to1$.} Then $\Wneu(E)=+\infty$; suppose $W_{C_0}(E)<+\infty$. There are real $M'$ and $R_0$ with $a(R)+b(R)\le M'$ for
$R\ge R_0$; in particular $F^{\square_R}<\infty$. By Corollary~\ref{cor:EB}, $F^{\square_R}\ge\frac1{4\pi}\int_{\square_R}|\Phi|^2-C_0^\ast n(R)$, hence
$\frac1{4\pi}\int_{\square_R}|\Phi|^2+(C_0-C_0^\ast)n(R)-C_0R^2\le M'R^2$ and, as $C_0\ge C_0^\ast$,
$\int_{\square_R}|\Phi|^2\le4\pi\max(M'+C_0,0)R^2$ for $R\ge R_0$. Proposition~\ref{prop:density} gives $n(R)/R^2\to1$, a contradiction.
\end{proof}

\begin{remark}
Theorem~\ref{thm:R2prime} uses only the results of \S\S2--3 (no screening). The constant $C_0$ of the Lectures (``the constant of (11.2.8)--(11.2.9)'')
may be any $C_0\ge\log2+\pi/32$; the Lean statement only asks for existence.
\end{remark}

%%%%%%%%%%%%%%%%%%%%%%%%%%%%%%%%%%%%%%%%%%%%%%%%%%%%%%%%%%%%%%%%%%%%%%%%%%%%%%%%%%%%%%%%%%%%%%%%%%%%%%%%%
\section{Screening of fields on a square}\label{sec:screening}

\subsection{Statements}
Throughout, $0<m_-\le m_+$ are fixed, $r_1=\sqrt2/4$, $R_1=r_1+\frac54$, $r_2=r_1+\frac12$, $R_2=r_2+\frac54$, $M_1=49$, $M_2=81$,
$a_j=a(r_j,m)=m\pi R_j^2$, $b_j=b(r_j)$ (Lemma~\ref{lem:localcount}), $r_0=\min\{\frac14,\frac1{16\sqrt{m_+}}\}$, and the constants
$\lambda_0,C_1,C_S$ are those of Definition~\ref{def:constants}.

\begin{theorem}[screening; \texttt{SandierSerfaty.Screening.exists\_screened\_neumannChargeField}]\label{thm:I4}
Put $\ell_0=8\lambda_0$, $C_{\rm scr}=512\,C_1$. Let $c\in\C$, $t,m,\lc,\lt>0$, $N\in\N$, $C\subset\C$, $\alpha:C\to\R$, $E$, and put
$Q=\sq(c,t)$, $\dep$ its depth, $U=\operatorname{int}\sq(c,t-1)=\{\dep>\frac12\}$. Assume
\begin{enumerate}[label=(H\arabic*)]
\item $m_-\le m\le m_+$, $mt^2=N$, $\ell_0\le\lc\le\lt$, $8\lt\le t$;
\item $C$ is locally finite and $E$ is a charge field on $U$ with density $m$ and charges $C$;
\item $0<\alpha_p\le\frac14$ for all $p\in C$, and $\alpha_p=\frac14$ if $p\notin\sq(c,t-2)$ (i.e.\ $\dep(p)<1$);
\item $\alpha_p+\alpha_q\le|p-q|$ for distinct $p,q\in C\cap\sq(c,t-2\lt)$ (i.e.\ $\dep\ge\lt$);
\item $E_\alpha\in L^2(Q)$;
\item $S:=\int_{\sq(c,t-\lt)\setminus\sq(c,t-4\lt)}|E_{1/4}|^2<\infty$ (the region is $\{\lt/2\le\dep<2\lt\}$) and $S\le\lc^3\lt/C_{\rm scr}$.
\end{enumerate}
Then there are a finite $Y\subset\operatorname{int}Q$ with $\#Y=N$, radii $\rho:Y\to(0,\frac14]$ with $\bar B(y,\rho_y)\subset\operatorname{int}Q$ and
$\rho_y+\rho_{y'}\le|y-y'|$ ($y\ne y'$), and a Neumann charge field $E'$ on $Q$ with density $m$ and charges $Y$, such that
\begin{equation}\label{eq:I4}
G_Q(\rho)[E';Y]\ \le\ G_Q(\alpha)[E;C]+C_S\Bigl(\frac{\lc S}{\lt}+t\lt+n_\partial\Bigr),\qquad n_\partial=\#\bigl(C\cap Q\setminus\sq(c,t-4\lt)\bigr).
\end{equation}
\end{theorem}

\begin{theorem}[internal form]\label{thm:I4int}
The conclusion of Theorem~\ref{thm:I4} holds if (H1) and (H6) are replaced by
\begin{enumerate}[label=(H\arabic*$'$)]
\item $m_-\le m\le m_+$, $mt^2=N$, $\lambda_0\le\lc$, $8\lc\le\lt$, $8\lt\le t$;
\setcounter{enumi}{5}
\item $S<\infty$ and $S\le\lc^3\lt/C_1$.
\end{enumerate}
\end{theorem}

\begin{proof}[Proof of Theorem~\ref{thm:I4} from Theorem~\ref{thm:I4int}]
Apply Theorem~\ref{thm:I4int} with $\lc'=\lc/8$: $\lc'\ge\lambda_0$, $8\lc'=\lc\le\lt$, $S\le\lc^3\lt/(512C_1)=\lc'^3\lt/C_1$, and the
error $\lc'S/\lt\le\lc S/\lt$.
\end{proof}

\begin{remark}[form of the hypotheses]
(a) The constants are fixed before the data. (b) Disjointness (H4) is required only in the kept zone
$\{\dep\ge\lt\}$; for the jellium lower bound the $\tilde r$-balls may overlap in the layer $\{1<\dep<2\}$ (the situation of the counterexample to Lemma~11.2 of the Lectures, see the errata record). (c) In Lean the hypotheses
are bundled in a structure, and the depth sets are the \texttt{sq} sets of \eqref{eq:depth-sets}.
\end{remark}

\subsection{Toolbox}

\begin{lemma}[rectangle solver]\label{lem:T2}
Let $H=\{x_0+s+i(y_0+r):0\le s\le a,\ 0\le r\le b\}$ (or its image under an affine isometry $x=z_0+M(s,r)$ with $M$ orthogonal, the field being pushed forward as $M\cdot(A,B)$;
the identity (b) and the bounds (a), (c) are invariant, and $M$ may be orientation-reversing, as for the charts $\Psi_j$ of Definition~\ref{def:pinwheel}) and $f\in L^1(H)$ real with $\int_Hf=0$. In the coordinates $(s,r)$ put $F_1(s)=\int_0^bf(s,r)\dd r$, $G(s)=\int_0^b|f(s,r)|\dd r$,
\[
A(s)=\frac1b\int_0^sF_1,\qquad B(s,r)=\int_0^r\Bigl(f(s,\varrho)-\frac{F_1(s)}b\Bigr)\dd\varrho,\qquad T2_H[f]:=(A,B)\mathbf 1_H .
\]
\begin{enumerate}[label=(\alph*)]
\item $A(0)=A(a)=0$, $B(s,0)=B(s,b)=0$; $|A|\le b^{-1}\int_0^aG$, $|B(s,r)|\le2G(s)$, and
  \[
  \int_H|T2_H[f]|^2\le\frac ab\Bigl(\int_0^aG\Bigr)^2+4b\int_0^aG^2\le\Bigl(\frac{a^2}b+4b\Bigr)\int_0^aG^2 .
  \]
\item For every Lipschitz $\varphi$ on $\C$ (compact support not needed): $\int_\C T2_H[f]\cdot\nabla\varphi=-\int_Hf\varphi$, i.e.\
  $\Div T2_H[f]=f\mathbf 1_H$ in $\mathcal D'(\C)$ (zero normal flux).
\item If $F_1=0$ a.e., then $A\equiv0$ and $|B(s,r)|\le\int_0^r|f(s,\cdot)|$.
\end{enumerate}
\end{lemma}
\begin{proof}
(a) $A(a)=b^{-1}\int_Hf=0$; $B(s,b)=F_1(s)-F_1(s)=0$; the bounds are immediate and
$\int_HA^2\le ab\cdot b^{-2}(\int G)^2$, $\int_HB^2\le b\int4G^2$; Cauchy--Schwarz gives the last form.
(b) For smooth $\varphi$, by Fubini on $\{0\le\varsigma\le s\le a\}$ and the fundamental theorem of calculus \emph{applied to $\varphi$ only}:
$\int_HA\,\partial_s\varphi=\frac1b\int_0^b\int_0^aF_1(\varsigma)[\varphi(a,r)-\varphi(\varsigma,r)]\dd\varsigma\dd r=-\frac1b\int_HF_1(s)\varphi$ (as $\int F_1=0$), and
$\int_HB\,\partial_r\varphi=\int_0^a\int_0^b(f-F_1/b)(s,\varrho)[\varphi(s,b)-\varphi(s,\varrho)]\dd\varrho\dd s=-\int_H(f-F_1/b)\varphi$ (as $\int_0^b(f-F_1/b)\dd\varrho=0$).
Adding gives the claim; all integrands are in $L^1$. For Lipschitz $\varphi$ use Lemma~\ref{lem:moll} ($U=\C$, $Y=T2_H[f]\in L^1$, $\nu=f\mathbf 1_H\dd x$); for
$\varphi$ without compact support multiply by a cutoff equal to $1$ near $H$. For the image under $\Phi(s,r)=z_0+M(s,r)$ the claims follow from the case $M=\mathrm{id}$
applied to $\varphi\circ\Phi$, since $\nabla(\varphi\circ\Phi)=M^{\mathsf T}(\nabla\varphi)\circ\Phi$, so $M(A,B)\cdot\nabla\varphi=(A,B)\cdot\nabla(\varphi\circ\Phi)$, and $|\det M|=1$.
(c) Immediate.
\end{proof}

\begin{lemma}[radial fields]\label{lem:radial}
For $r>0$ let $u^r=(\pi r^2)^{-1}\mathbf 1_{B(0,r)}$, $W^r(y)=yr^{-2}\mathbf 1_{B(0,r)}(y)$ and $K^r=\nabla f_r+W^r$, i.e.\ $K^r(y)=-y|y|^{-2}(1-|y|^2/r^2)_+$.
For every Lipschitz $\varphi$ with compact support:
\begin{enumerate}[label=(\alph*)]
\item $\int W^r\cdot\nabla\varphi=2\pi\bigl(\avg_{\partial B(0,r)}\varphi-\int u^r\varphi\bigr)$;
\item $\int K^r\cdot\nabla\varphi=2\pi\bigl(\varphi(0)-\int u^r\varphi\bigr)$, i.e.\ $-\Div K^r=2\pi(\delta_0-u^r)$;
\item $|W^r|\le1/r$, $\int|W^r|^2=\pi/2$, $\int|W^r|=2\pi r/3$; $K^r-\nabla f_r=W^r$ and $\supp K^r\subset\bar B(0,r)$.
\end{enumerate}
For $r=\frac14$ write $W=W^{1/4}=16y\mathbf 1_{B(0,1/4)}$, $u=u^{1/4}=\frac{16}\pi\mathbf 1_{B(0,1/4)}$, $K=K^{1/4}$: $|W|\le4$, $\int|W|=\pi/6$.
\end{lemma}
\begin{proof}
(a) For smooth $\varphi$ apply the divergence theorem to $\varphi\,y/r^2$ on $B(0,r)$: $\Div(y/r^2)=2/r^2$ and $y/r^2\cdot\nu=1/r$ on the circle, so
$\int_{B_r}r^{-2}y\cdot\nabla\varphi=\frac1r\int_{\partial B_r}\varphi-\frac2{r^2}\int_{B_r}\varphi$; then Lemma~\ref{lem:moll}. (b) Add Lemma~\ref{lem:f}(c).
(c) Polar coordinates.
\end{proof}

\begin{lemma}[averaging]\label{lem:avg}
If $h\ge0$ is measurable on $[s_1,s_2]$ and $\int_{s_1}^{s_2}h\le M$, then $h(s)\le M/(s_2-s_1)$ for some $s\in[s_1,s_2]$.
\end{lemma}
\begin{proof}Otherwise $\int h>M$.\end{proof}

\subsection{Standing data; good cut; keep and drop; soft cut}
Until the end of \S\ref{sec:screening} fix data satisfying the hypotheses of Theorem~\ref{thm:I4int}, and write $\Phi=E_{1/4}$ (all of $C$
truncated at $\frac14$), $\sigma:=S/\lt$. Definition~\ref{def:constants} gives $\lc\ge\lambda_0\ge\max\{36,8/m,2/\sqrt m\}$, hence
$\lt\ge8\lc\ge288$, $t\ge8\lt$.

\begin{lemma}[good cut]\label{lem:goodcut}
There is $d\in[\lt,\frac32\lt]$ such that the ring $R_d:=\{x:d-5\le\dep(x)\le d+11\}$ satisfies $R_d\subset\{\lt/2\le\dep<2\lt\}\subset U$ and
\[
S_d:=\int_{R_d}|\Phi|^2\le32\,\sigma .
\]
\end{lemma}
\begin{proof}
$\lt-5\ge\lt/2$ and $\frac32\lt+11<2\lt$ as $\lt\ge288$. By Tonelli,
$\int_{\lt}^{3\lt/2}S_d\dd d=\int|\Phi(x)|^2\,\bigl|\{d\in[\lt,\tfrac32\lt]:|\dep(x)-d-3|\le8\}\bigr|\dd x\le16S$; apply Lemma~\ref{lem:avg}.
\end{proof}

Fix this $d$ and name the regions
\[
\Omega_+=\{\dep>d\},\quad \Omega_{\rm c}=\{\dep>d+8\},\quad A=\{d-1<\dep\le d\},\quad O_-=Q\cap\{\dep\le d-1\},
\]
\[
F=Q\cap\{\dep\le d+8\}.
\]

\begin{definition}[keep and drop]\label{def:keep}
$K=\{p\in C:\dep(p)\ge d+6\}$, $D^\ast=\{q\in C:d-\frac54<\dep(q)<d+6\}$; $\beta_p=\alpha_p$ for $p\in K$, $\beta_p=\frac14$ for $p\in C\setminus K$;
$F_p=f_{\alpha_p,1/4}(\cdot-p)$ for $p\in K$.
\end{definition}

\begin{lemma}[keep and drop]\label{lem:keep}
\begin{enumerate}[label=(\alph*)]
\item $K\subset\operatorname{int}Q$ is finite; $\alpha_p+\alpha_q\le|p-q|$ for distinct $p,q\in K$; $\bar B(p,\alpha_p)\subset\{\dep\ge d+\frac{23}4\}$ for $p\in K$.
\item $D^\ast\cup\bigl((C\setminus K)\cap Q\bigr)\subset Q\cap\{\dep<2\lt\}$; in particular $\#((C\setminus K)\cap Q)\le n_\partial$.
\item $E_\beta\in L^2(Q)$ and $\frac12\int_Q|E_\beta|^2+\sum_{p\in K}\tau_m(\alpha_p)\le G_Q(\alpha)[E]+c_m\,n_\partial$.
\item $E_\beta=\Phi-\sum_{p\in K}\nabla F_p$ on $\C$.
\item Every $q\in C\setminus(K\cup D^\ast)$ has $\dep(q)\le d-\frac54$, so $B(q,\frac14)\cap\{\dep>d-1\}=\emptyset$.
\end{enumerate}
\end{lemma}
\begin{proof}
(a) $d+6\ge\lt$, so (H4) applies; $\dep$ is $1$-Lipschitz. (b) $d+6\le\frac32\lt+6<2\lt$, and $d-\frac54>0$.
(c) Proposition~\ref{prop:mono} with $D=Q$, $U$, $\alpha\le\beta$: the changed points $p\in C\setminus K$ with $\alpha_p<\frac14$ have $\dep(p)\ge1$ by (H3), so
$\bar B(p,\frac14)\subset\{\dep\ge\frac34\}\subset U\cap Q$; (H5) gives $E_\alpha\in L^2(Q)$. By Corollary~\ref{cor:mono}(a), $G_Q(\beta)\le G_Q(\alpha)$, and
$G_Q(\beta)=\frac12\int_Q|E_\beta|^2+\sum_K\tau_m(\alpha_p)-c_m\#((C\setminus K)\cap Q)$; use (b).
(d) $f_{1/4}-f_{\alpha_p}=-f_{\alpha_p,1/4}$. (e) Definitions.
\end{proof}

\begin{definition}[soft cut]\label{def:soft}
$\psi=\min\{1,\max\{0,\dep-d+1\}\}$ ($1$-Lipschitz, $=1$ on $\{\dep\ge d\}$, $=0$ on $\{\dep\le d-1\}$, $\nabla\psi=0$ a.e.\ off $A$,
$\supp\psi=\{\dep\ge d-1\}$ compact in $U$). With $u_q=u(\cdot-q)$, $W_q=W(\cdot-q)$, $K_q=K(\cdot-q)$ (Lemma~\ref{lem:radial}):
\[
X=E-\sum_{q\in D^\ast}K_q\ \text{on }U,\qquad g=2\pi m(1-\psi)+2\pi\psi\sum_{q\in D^\ast}u_q-\nabla\psi\cdot X\ \text{on }Q,
\]
($\nabla\psi\cdot X:=0$ where $\nabla\psi=0$), $n(x)=\#\{q\in D^\ast:|x-q|<\frac14\}$, $N_2=\int_\C n^2$, $X_A=\int_A|X|^2$.
\end{definition}

\begin{lemma}[soft cut]\label{lem:soft}
\begin{enumerate}[label=(\alph*)]
\item On $\{\dep>d-1\}$, a.e.: $X-\sum_{p\in K}\nabla f_{\alpha_p}(\cdot-p)=E_\beta-\sum_{q\in D^\ast}W_q$. On $A$: $X=\Phi-\sum_{D^\ast}W_q$ and $|X|\le|\Phi|+4n$.
\item For every Lipschitz $\varphi$ on $\C$: $\int\psi X\cdot\nabla\varphi=2\pi\sum_{p\in K}\varphi(p)+\int_Q(g-2\pi m)\varphi$.
\item $g\in L^2(Q)$; $g=2\pi m$ on $O_-$; $g=2\pi\sum_{D^\ast}u_q$ on $Q\cap\{\dep\ge d\}$; $g=0$ on $\Omega_{\rm c}$;
  $|g|\le2\pi m+32n+|\nabla\psi\cdot X|$.
\item $\int_Qg=2\pi(N-\#K)$.
\end{enumerate}
\end{lemma}
\begin{proof}
(a) On $\{\dep>d-1\}$, $E_\beta=E-\sum_K\nabla f_{\alpha_p}(\cdot-p)-\sum_{D^\ast}\nabla f_{1/4}(\cdot-q)$ by Lemma~\ref{lem:keep}(e), and $K_q=\nabla f_{1/4}(\cdot-q)+W_q$.
On $A$ no kept disc is present (Lemma~\ref{lem:keep}(a)) and $\Phi=E_\beta$ there; $|\sum W_q|\le4n$.
(b) $\psi\varphi$ is a Lipschitz test function on $U$ (cut off $\varphi$ outside a neighbourhood of $\supp\psi$). By \eqref{eq:pair0} and Lemma~\ref{lem:radial}(b),
$\int X\cdot\nabla(\psi\varphi)=2\pi[\sum_C\psi\varphi(p)-m\int\psi\varphi]-2\pi\sum_{D^\ast}[\psi\varphi(q)-\int u_q\psi\varphi]
=2\pi[\sum_K\varphi(p)+\sum_{D^\ast}\int u_q\psi\varphi-m\int\psi\varphi]$, since $\psi=1$ on $K$ and $\psi=0$ on $C\setminus(K\cup D^\ast)$. Insert
$\nabla(\psi\varphi)=\psi\nabla\varphi+\varphi\nabla\psi$ a.e.
(c) $X\in L^2(A)$ by (a) ($A\subset R_d$, $W$ bounded, $n$ bounded on compacts); the formulas follow from the properties of $\psi$ and
$\supp u_q\subset\{\dep<d+\frac{25}4\}$; $2\pi\cdot\frac{16}\pi=32$.
(d) Take $\varphi\in C^\infty_c$ equal to $1$ near $Q$ in (b): $\nabla\varphi=0$ on $\supp\psi X$, so $0=2\pi\#K+\int_Qg-2\pi mt^2$.
\end{proof}

\begin{lemma}[ring estimates]\label{lem:ring}
Let $\mathcal G=\{x\in\Z^2/2:d-\frac32<\dep(x)<d+\frac{25}4\}$ and $k_j(x)=\#(C\cap\bar B(x,r_j))$.
\begin{enumerate}[label=(\alph*)]
\item $\#\mathcal G\le132t$; for $x\in\mathcal G$: $\bar B(x,R_2)\subset R_d\cap U$ and $k_j(x)\le a_j+b_j\|\Phi\|_{L^2(B(x,R_j))}$; every $q\in D^\ast$
  has some $x(q)\in\mathcal G$ with $|q-x(q)|_\infty\le\frac14$.
\item $N_2\le\frac\pi{16}\bigl(264a_2^2t+2b_2^2M_2S_d\bigr)$.
\item $X_A\le2S_d+32N_2$.
\item For every closed axis-parallel rectangle $T$ with sides $w\ge9$ and $9$:
  $n_D(T):=\#\{q\in D^\ast:B(q,\frac14)\cap T\ne\emptyset\}\le21(2w+3)a_1+b_1\sqrt{21(2w+3)M_1S_d}$ and
  $\int_{T\cap A}|X|\le\sqrt{|T\cap A|\,S_d}+\frac\pi6n_D(T)$.
\end{enumerate}
\end{lemma}
\begin{proof}
(a) Lemma~\ref{lem:grid}(d) with $s_2-s_1=\frac{31}4$; $\bar B(x,R_2)\subset\{d-3.61<\dep<d+8.36\}\subset R_d$ ($R_2<2.11$); Lemma~\ref{lem:localcount}
($\Phi\in L^2(R_d)$ by Lemma~\ref{lem:goodcut}); Lemma~\ref{lem:grid}(a) and $\dep(x)\in[\dep(q)-\frac14,\dep(q)+\frac14]$.
(b) $N_2=\sum_{q,q'\in D^\ast}|B(q,\frac14)\cap B(q',\frac14)|\le\frac\pi{16}\#\{(q,q'):|q-q'|<\frac12\}$. Grouping $q$ by $x(q)$: at most $k_1(x)\le k_2(x)$
points $q$ have $x(q)=x$, and each $q'$ with $|q'-q|<\frac12$ lies in $\bar B(x,r_2)$. So $N_2\le\frac\pi{16}\sum_{\mathcal G}k_2^2\le\frac\pi{16}\sum_{\mathcal G}(2a_2^2+2b_2^2\|\Phi\|^2_{L^2(B(x,R_2))})$;
use (a) and Lemma~\ref{lem:grid}(b).
(c) Lemma~\ref{lem:soft}(a), $A\subset R_d$. (d) For such $q$, $|q-y|_\infty<\frac14$ for some $y\in T$, so $x(q)\in T':=T+[-\frac12,\frac12]^2$ (sides
$w+1$, $10$), and $\#(\mathcal G\cap T')\le21(2w+3)$ (Lemma~\ref{lem:grid}(c)). Sum (a) over $\mathcal G\cap T'$ with Cauchy--Schwarz and Lemma~\ref{lem:grid}(b).
The second bound: (a) of Lemma~\ref{lem:soft}, Cauchy--Schwarz, $\int|W_q|=\pi/6$.
\end{proof}

\subsection{The frame, columns and carries}

\begin{definition}[pinwheel]\label{def:pinwheel}
Let $a_\square=\frac t2-d-8$ ($>0$ since $t\ge8\lt$, $d\le\frac32\lt$), $L=t-d-8$, $w_f=d+8$, $\omega=-i$. For $j\in\{0,1,2,3\}$ let
$\Psi_j(u,v)=c+\omega^j\bigl((u-a_\square)+i(\tfrac t2-v)\bigr)$ (an isometry $\R^2\to\C$) and $R_j=\Psi_j([0,L]\times[0,w_f])$.
Top zone $T^{(j)}=\Psi_j([0,L]\times[d-1,w_f])$, $T=\bigcup_jT^{(j)}$.
\end{definition}

\begin{lemma}[pinwheel geometry]\label{lem:pinwheel}
\begin{enumerate}[label=(\alph*)]
\item $R_0=c+[-a_\square,\frac t2]\times[a_\square,\frac t2]$ (as $x+iy$), $R_{j+1}=c+\omega(R_j-c)$; the $R_j$ have pairwise disjoint interiors and
  $\bigcup_jR_j=F=Q\setminus\operatorname{int}\sq(c,2a_\square)$.
\item $\dep(\Psi_j(u,v))=\min\{v,L-u\}$ on $[0,L]\times[0,w_f]$.
\item Hence, in the coordinates of $R_j$: $\{\dep\ge d-1\}\cap R_j=\Psi_j([0,L-d+1]\times[d-1,w_f])\subset T^{(j)}$;
  $(A\cup\Omega_+)\cap R_j\subset T^{(j)}$; $\Psi_j([0,L]\times[0,d-1])\subset O_-$; $|T|\le36t$.
\item For $0\le w\le w_f$: $\Psi_j([L-w,L]\times[w_f-1,w_f])\cup\Psi_{j+1}([0,1]\times[0,w])$ is a closed rectangle with sides $w$ and $2$, bisected
  by the common segment $\Psi_j([L-w,L]\times\{w_f\})=\Psi_{j+1}(\{0\}\times[0,w])$ (indices mod $4$).
\end{enumerate}
\end{lemma}
\begin{proof}
(a) For $j=0$: $x=u-a_\square\in[-a_\square,\frac t2]$, $y=\frac t2-v\in[a_\square,\frac t2]$. Multiplication by $\omega$ maps $x+iy$ to $y-ix$, so
$R_1=c+[a_\square,\frac t2]\times[-\frac t2,a_\square]$, $R_2=c+[-\frac t2,a_\square]\times[-\frac t2,-a_\square]$, $R_3=c+[-\frac t2,-a_\square]\times[-a_\square,\frac t2]$. A point of
$Q$ with $|x-c|_\infty\ge a_\square$ lies in $R_0$ if $y\ge a_\square$, $x\ge-a_\square$; in $R_3$ if $y\ge-a_\square$, $x\le-a_\square$; in $R_1$ if $x\ge a_\square$, $y\le a_\square$;
in $R_2$ if $y\le-a_\square$, $x\le a_\square$ (coordinates relative to $c$); overlaps are segments. $F=\{\dep\le d+8\}\cap Q=Q\setminus\operatorname{int}\sq(c,t-2d-16)$.
(b) For $j=0$, $|y|=\frac t2-v$ and $|x|=|u-a_\square|$, so $\dep=\min\{v,\frac t2-|u-a_\square|\}$; if $u\le2a_\square$ then $\frac t2-|u-a_\square|\ge\frac t2-a_\square=w_f\ge v$, and if
$u\ge2a_\square$ then $\frac t2-|u-a_\square|=L-u$; in both cases $\dep=\min\{v,L-u\}$ since $L-u\ge w_f\ge v$ for $u\le2a_\square$. Rotations about $c$ by multiples of
$\frac\pi2$ preserve $|\cdot-c|_\infty$. (c) From (b): $\dep\ge d-1$ iff $v\ge d-1$ and $L-u\ge d-1$; $A,\Omega_+\subset\{\dep>d-1\}$; $|T|\le4\cdot9L\le36t$.
(d) For $j=0$ the two sets are $c+[\frac t2-w,\frac t2]\times[a_\square,a_\square+1]$ and $c+[\frac t2-w,\frac t2]\times[a_\square-1,a_\square]$
($\Psi_1(u,v)=c+(\frac t2-v)+i(a_\square-u)$); rotate.
\end{proof}

\begin{definition}[columns]\label{def:columns}
A column partition of $R_j$ is $0=u_0<u_1<\dots<u_{K_j}=L$; for $I_k=[u_{k-1},u_k]$, $w_k=|I_k|$:
$\mathrm{Col}_k=\Psi_j(I_k\times[0,w_f])$, top part $T_k=\Psi_j(I_k\times[d-1,w_f])$, filler part $P_k=\Psi_j(I_k\times[0,d-1])$,
compensation strip $H'_k=\Psi_j(I_k\times[d-1-\lc,d-1])$, flux cell $H_k=H'_k\cup T_k$. For an interval $I\subset[0,L]$, $T_I=\Psi_j(I\times[d-1,w_f])$.
\end{definition}

\begin{definition}[conditions on the parameters]\label{def:K}
\begin{enumerate}[label=(K\arabic*)]
\item for every $w\in[\lc,3\lc]$, with $n_D^{\max}(w):=21(2w+3)a_1+b_1\sqrt{21(2w+3)M_1\cdot32\lc^3/C_1}$:
\[
\Theta(w):=18\pi mw+\tfrac{13\pi}6\,n_D^{\max}(w)+\sqrt{64w\lc^3/C_1}+2\pi\ \le\ \tfrac\pi2\,mw\lc ;
\]
\item $1/\lc\le m/8$; \quad (K3) $\lc\sqrt m\ge2$; \quad (K4) $\lc\ge36$.
\end{enumerate}
\end{definition}

\begin{lemma}[strip estimate]\label{lem:strip}
Assume (K1). Let $I\subset[0,L]$ be an interval of length $w\in[\lc,3\lc]$ and $h=g-\kappa\mathbf 1_{A^\flat}$ on $T_I$ with $A^\flat\subset T_I$, $\kappa\ge0$,
$\kappa|A^\flat|\le2\pi$. Then $\int_{T_I}|h|\le\Theta(w)\le\frac\pi2mw\lc$.
\end{lemma}
\begin{proof}
$T_I$ is a $w\times9$ rectangle. On $T_I$, $|g|\le2\pi m(1-\psi)+2\pi\psi\sum u_q+|\nabla\psi\cdot X|$ with $\int_{T_I}2\pi m(1-\psi)\le18\pi mw$,
$\int_{T_I}2\pi\sum u_q\le2\pi n_D(T_I)$ ($\int u_q=1$), and $\int_{T_I}|\nabla\psi\cdot X|\le\int_{T_I\cap A}|X|\le\sqrt{2wS_d}+\frac\pi6n_D(T_I)$
by Lemma~\ref{lem:ring}(d), since $|T_I\cap A|\le w+9\le2w$ (Lemma~\ref{lem:pinwheel}(b): $A\cap R_j\subset\{d-1<v\le d\}\cup\{d-1<L-u\le d\}$). Finally
$n_D(T_I)\le n_D^{\max}(w)$ and $S_d\le32\sigma\le32\lc^3/C_1$ by Lemma~\ref{lem:goodcut} and (H6$'$).
\end{proof}

\begin{lemma}[column cuts with carries]\label{lem:cuts}
Assume (K1),(K2),(K4). There are column partitions of $R_0,\dots,R_3$ and $\theta_0,\theta_1,\theta_2\in[0,1)$ such that, with $w'_j$ the width of the last column of
$R_j$, $A_j=\Psi_j([L-w'_j,L]\times[w_f-1,w_f])$, $A'_j=\Psi_{j+1}([0,1]\times[0,w'_j])$ and
\[
\tilde g:=g-\sum_{j=0}^2\frac{2\pi\theta_j}{w'_j}\bigl(\mathbf 1_{A_j}-\mathbf 1_{A'_j}\bigr),
\]
\begin{enumerate}[label=(\alph*)]
\item all columns have width in $[\lc,2\lc]$, except the last column of each $R_j$, whose width lies in $[\lc,3\lc)$;
\item $n_k:=\frac1{2\pi}\int_{\mathrm{Col}_k}\tilde g\in\N$ for every column of every $R_j$, and $\sum n_k=N-\#K$;
\item $A_j\subset T_{\rm last}(R_j)\cap O_-$, $A'_j\subset(P_{\rm first}(R_{j+1})\setminus H'_{\rm first})\cap O_-$; $\tilde g=2\pi m$ on every $P_k\setminus\bigcup_jA'_j$ and
  $\tilde g=2\pi m+2\pi\theta_j/w'_j\in[2\pi m,2\pi m+2\pi/\lc]$ on $A'_j$;
\item for every interval $I\subset[0,L]$ of length $w\in[\lc,3\lc]$: $|\int_{T_I}\tilde g|\le\frac\pi2mw\lc$ and $\frac1{2\pi}\int_{\Psi_j(I\times[0,w_f])}\tilde g\ge6mw\lc$.
\end{enumerate}
\end{lemma}
\begin{proof}
(c) and the first part of (d) do not depend on the choices: $A_j\subset\{\dep\le L-u\le w'_j\}$ with $w'_j<3\lc<d-1-\lc$ (as $d\ge8\lc$), $A'_j\subset\{v\le w'_j\}$;
on $T_I$, $\tilde g=g-\kappa\mathbf 1_{A_j\cap T_I}$ with $\kappa|A_j|\le2\pi$, so Lemma~\ref{lem:strip} applies. For the mass: on $\Psi_j(I\times[0,d-1])\subset O_-$,
$\tilde g\ge2\pi m$; hence $\frac1{2\pi}\int\tilde g\ge mw(d-1)-\frac14mw\lc\ge mw(7\lc-1-\frac14\lc)\ge6mw\lc$.

Construction, for $j=0,1,2,3$ in turn. Let $\tilde g^{(j)}=g+\frac{2\pi\theta_{j-1}}{w'_{j-1}}\mathbf 1_{A'_{j-1}}$ on $R_j$ ($\theta_{-1}:=0$), already determined, and
$\mu(u)=\frac1{2\pi}\int_{\Psi_j([0,u]\times[0,w_f])}\tilde g^{(j)}$, continuous in $u$ ($\tilde g^{(j)}\in L^1$). Put $u_0=0$. While $L-u_{k-1}\ge3\lc$: on
$[u_{k-1}+\lc,u_{k-1}+2\lc]$ the continuous function $\mu(\cdot)-\mu(u_{k-1})$ increases between the end points by the mass of a strip of width $\lc$, which is
$\ge6m\lc^2\ge1$ by (d) applied to $\tilde g^{(j)}$ (Lemma~\ref{lem:strip} with $\kappa=0$ bounds the top part of $\tilde g^{(j)}$, and $\tilde g^{(j)}\ge2\pi m$ on $\Psi_j(I\times[0,d-1])$; the non-last columns do not meet
$A_j$, so their $\tilde g^{(j)}$- and $\tilde g$-masses coincide). By the intermediate value theorem
choose $u_k$ there with $\mu(u_k)-\mu(u_{k-1})\in\Z$; it is $\ge6mw_k\lc>0$. When $L-u_{k-1}<3\lc$, the last column is $[u_{k-1},L]$; its width is $\ge\lc$
(the previous step left $L-u_{k-1}\ge3\lc-2\lc$; and $L\ge3\lc$). Its mass $M_{\rm last}=\mu(L)-\mu(u_{k-1})\ge6m\lc^2$. For $j\le2$ put
$\theta_j=M_{\rm last}-\lfloor M_{\rm last}\rfloor$; subtracting $\frac{2\pi\theta_j}{w'_j}\mathbf 1_{A_j}$ makes the last column's mass $\lfloor M_{\rm last}\rfloor\in\N$.
$A'_j$ lies in the first column of $R_{j+1}$ because $u_1\ge\lc\ge1$. For $j=3$: $\frac1{2\pi}\int_Q\tilde g=\frac1{2\pi}\int_Qg=N-\#K\in\Z$
(Lemma~\ref{lem:soft}(d); the carries have zero total as $|A_j|=|A'_j|=w'_j$, and $\tilde g=g=0$ on $\Omega_{\rm c}$), and all other columns have integer
mass, so the last column of $R_3$ has integer mass; it is positive. This proves (a), (b).
\end{proof}

\subsection{Flux cells, carries, filler}

\begin{lemma}[flux cells]\label{lem:flux}
For every column (of every $R_j$) put $\mu_k=m+\frac1{2\pi w_k\lc}\int_{T_k}\tilde g$, $f_k=\tilde g$ on $T_k$, $f_k=2\pi(m-\mu_k)$ on $H'_k$, and
$V_k=T2_{H_k}[f_k]$ in the coordinates $\Psi_j$ ($s$-axis along $u$). Then
\begin{enumerate}[label=(\alph*)]
\item $|\mu_k-m|\le m/4$; $\int_{H_k}f_k=0$ and $\int V_k\cdot\nabla\varphi=-\int_{H_k}f_k\varphi$ for Lipschitz $\varphi$;
\item $\|V_k\|^2_{L^2}\le504\,\lc\int_{T_k}|\tilde g|^2$, hence $V_{\rm flux}:=\sum_k\|V_k\|^2\le504\,\lc\,\hat G$ with $\hat G:=\int_T|\tilde g|^2$;
\item $\hat G\le4\bigl(144\pi^2m^2t+1024N_2+X_A+12\pi^2/\lc\bigr)$.
\end{enumerate}
\end{lemma}
\begin{proof}
(a) Lemma~\ref{lem:cuts}(d) with $I=I_k$; $\int_{H_k}f_k=\int_{T_k}\tilde g+2\pi(m-\mu_k)w_k\lc=0$; Lemma~\ref{lem:T2}(b).
(b) With $L(u)=\int_{d-1}^{w_f}|\tilde g(\Psi_j(u,v))|\dd v$: $G_k(u)=L(u)+2\pi|m-\mu_k|\lc$ and $2\pi|m-\mu_k|\lc=|\int_{T_k}\tilde g|/w_k\le\bar L_k$ (mean of $L$ on
$I_k$), so $\int_{I_k}G_k^2\le2\int L^2+2w_k\bar L_k^2\le4\int_{I_k}L^2\le36\int_{T_k}|\tilde g|^2$. Lemma~\ref{lem:T2}(a) with $a=w_k\le3\lc$, $b=\lc+9$:
$a^2/b+4b\le9\lc+4\lc+36\le14\lc$ by (K4). The top parts have disjoint interiors.
(c) On $T$, $\tilde g=2\pi m(1-\psi)+2\pi\psi\sum u_q-\nabla\psi\cdot X-\sum_j\frac{2\pi\theta_j}{w'_j}\mathbf 1_{A_j}$; $(\sum_1^4x_i)^2\le4\sum x_i^2$; $|T|\le36t$;
$(2\pi\cdot\frac{16}\pi)^2=1024$; $\int_T|\nabla\psi\cdot X|^2\le X_A$; $\sum_j(2\pi/w'_j)^2|A_j|\le12\pi^2/\lc$.
\end{proof}

\begin{lemma}[carries]\label{lem:carry}
For $j\in\{0,1,2\}$ let $H^c_j=A_j\cup A'_j$ (a $w'_j\times2$ rectangle, Lemma~\ref{lem:pinwheel}(d)), $f^c_j=\frac{2\pi\theta_j}{w'_j}(\mathbf 1_{A_j}-\mathbf 1_{A'_j})$ and
$V^c_j=T2_{H^c_j}[f^c_j]$ with the $s$-axis along the common segment. Then $F_1\equiv0$, $|V^c_j|\le2\pi\theta_j/w'_j$,
$\int V^c_j\cdot\nabla\varphi=-\int f^c_j\varphi$ for Lipschitz $\varphi$, $\|V^c_j\|^2=\frac{8\pi^2\theta_j^2}{3w'_j}\le\frac{8\pi^2}{3\lc}$, and
$H^c_j\subset O_-\cap\{\dep\le3\lc\}$; the three carry cells are pairwise disjoint.
\end{lemma}
\begin{proof}
Every segment of $H^c_j$ across the common segment meets $A_j$ and $A'_j$ in unit lengths with opposite constant values, so $F_1=0$; Lemma~\ref{lem:T2}(c):
$B$ is the tent $r\mapsto-\frac{2\pi\theta_j}{w'_j}\min(r,2-r)$ (orientation $A'_j$ first), whose square integrates to $\frac23(2\pi\theta_j/w'_j)^2$ on $[0,2]$; multiply by
$w'_j$. Location: Lemma~\ref{lem:cuts}(c), $\dep\le w'_j\le3\lc$. Disjointness: in coordinates relative to $c$ (as $x+iy$),
$H^c_0\subset[\frac t2-3\lc,\frac t2]\times[a_\square-1,a_\square+1]$,
$H^c_1\subset[a_\square-1,a_\square+1]\times[-\frac t2,-\frac t2+3\lc]$,
$H^c_2\subset[-\frac t2,-\frac t2+3\lc]\times[-a_\square-1,-a_\square+1]$.
They are pairwise disjoint because $a_\square-1>3\lc-\frac t2$ and $\frac t2-3\lc>a_\square+1$, i.e.\ $d+8>3\lc+1$.
\end{proof}

\begin{lemma}[filler partition]\label{lem:fillpart}
Let $P=\Psi_j([u_-,u_-+w]\times[0,h])$, $\mu$ measurable on $P$ with $\frac m2\le\mu\le\frac{3m}2$, $n:=\int_P\mu\in\N$, $w\sqrt m\ge2$, $h\sqrt m\ge2$.
Then $P$ is the union of $n$ closed rectangles $P_1,\dots,P_n$ with sides parallel to those of $P$ and pairwise disjoint interiors, with $\int_{P_i}\mu=1$, sides along $u$
in $[\frac1{6\sqrt m},\frac6{\sqrt m}]$ and along $v$ in $[\frac1{3\sqrt m},\frac4{\sqrt m}]$.
\end{lemma}
\begin{proof}
$n_r:=\lfloor w\sqrt m\rfloor\in[\frac12w\sqrt m,w\sqrt m]$ (as $w\sqrt m\ge1$) and $n\ge\frac m2wh\ge w\sqrt m\ge n_r$. Let $K'=\lfloor n/n_r\rfloor\ge1$, $s=n-K'n_r\in[0,n_r)$,
and split $n=\sum_{i=1}^{K'}n_i$ with $n_i\in\{n_r+\lfloor s/K'\rfloor,n_r+\lfloor s/K'\rfloor+1\}$; then $n_r\le n_i\le2n_r$. The cumulative mass
$\mathsf R(v)=\int_{\Psi_j([u_-,u_-+w]\times[0,v])}\mu$ is continuous and strictly increasing from $0$ to $n$ ($\mu\ge m/2>0$); cut the rows at
$v_i=\mathsf R^{-1}(n_1+\dots+n_i)$. Row $i$ has mass $n_i$, so its height $h_i$ satisfies $\frac m2wh_i\le n_i\le\frac{3m}2wh_i$, i.e.\
$h_i\in[\frac{2n_r}{3mw},\frac{4n_r}{mw}]\subset[\frac1{3\sqrt m},\frac4{\sqrt m}]$. In row $i$ cut along $u$ at the points where the cumulative mass of the row equals
$1,2,\dots,n_i-1$ (again continuous and strictly increasing); a piece of mass $1$ and height $h_i$ has width in $[\frac2{3mh_i},\frac2{mh_i}]\subset[\frac1{6\sqrt m},\frac6{\sqrt m}]$.
\end{proof}

\begin{lemma}[piece fields]\label{lem:piece}
Let $P_i$ be a piece of Lemma~\ref{lem:fillpart} with centre $z$, $u_z=u^{r_0}(\cdot-z)$, and
$\Psi_{P_i}=K^{r_0}(\cdot-z)+T2_{P_i}[2\pi(\mu-u_z)]$. Then $\bar B(z,r_0)\subset\operatorname{int}P_i$; for Lipschitz $\varphi$,
$\int\Psi_{P_i}\cdot\nabla\varphi=2\pi(\varphi(z)-\int_{P_i}\mu\varphi)$; and
\[
\int|\Psi_{P_i}-\nabla f_{r_0}(\cdot-z)|^2\le2\Bigl(\frac\pi2+e_T(m)\Bigr),\qquad e_T(m)=1056\pi^2+\frac{256\pi}{r_0\sqrt m}.
\]
\end{lemma}
\begin{proof}
$r_0\le\frac1{16\sqrt{m_+}}\le\frac1{16\sqrt m}<\frac1{12\sqrt m}\le\frac12\min$(sides). The identity: Lemma~\ref{lem:radial}(b) and Lemma~\ref{lem:T2}(b)
($\int_{P_i}2\pi(\mu-u_z)=0$). Energy: $\Psi_{P_i}-\nabla f_{r_0}(\cdot-z)=W^{r_0}(\cdot-z)+T2_{P_i}[\cdot]$ and $\int|W^{r_0}|^2=\pi/2$. For the T2 part (sides $a'$ along
$u$, $b'$ along $v$): $G\le2\pi\cdot\frac{3m}2b'+2\pi\cdot\frac{2r_0}{\pi r_0^2}=3\pi mb'+\frac4{r_0}$, $\int G\le2\pi(1+1)=4\pi$, $a'/b'\le18$, $b'\le4/\sqrt m$; so
$\|T2\|^2\le18\cdot16\pi^2+4b'\cdot4\pi(12\pi\sqrt m+4/r_0)\le288\pi^2+768\pi^2+256\pi/(r_0\sqrt m)$.
\end{proof}

\subsection{Assembly}

\begin{definition}[the screened configuration]\label{def:assembly}
For every column, apply Lemma~\ref{lem:fillpart} to the filler part $P_k$ (sides $w_k$, $h=d-1$) with $\mu=\tilde g/2\pi$ on $P_k\setminus H'_k$ and
$\mu=\mu_k$ on $H'_k$. Let $\mathcal P$ be the set of all pieces, $Z$ their centres, $Y=K\cup Z$, $\rho_p=\alpha_p$ ($p\in K$), $\rho_z=r_0$ ($z\in Z$), and
\[
E'=\psi X+\sum_kV_k+\sum_{j=0}^2V^c_j+\sum_{P\in\mathcal P}\Psi_P .
\]
\end{definition}

\begin{lemma}[well-posedness and points]\label{lem:points}
Assume (K1)--(K4). The hypotheses of Lemma~\ref{lem:fillpart} hold for every filler part; $\#Z=\sum n_k=N-\#K\le9mt\lt$; $\#Y=N$; $Y\subset\operatorname{int}Q$;
$0<\rho\le\frac14$; $\bar B(y,\rho_y)\subset\operatorname{int}Q$; $\rho_y+\rho_{y'}\le|y-y'|$ for $y\ne y'$.
\end{lemma}
\begin{proof}
Density: $\mu_k\in[\frac34m,\frac54m]$ (Lemma~\ref{lem:flux}(a)), $\mu\in[m,m+1/\lc]\subset[m,\frac98m]$ elsewhere (Lemma~\ref{lem:cuts}(c), (K2)).
Mass: as $\tilde g=2\pi m$ on $H'_k$, $\int_{P_k}\mu=\frac1{2\pi}\int_{P_k\setminus H'_k}\tilde g+\mu_kw_k\lc=\frac1{2\pi}\bigl(\int_{P_k\setminus H'_k}\tilde g+2\pi mw_k\lc+\int_{T_k}\tilde g\bigr)=n_k\in\N$.
Sides: $w_k\sqrt m\ge\lc\sqrt m\ge2$ (K3), $h\sqrt m=(d-1)\sqrt m\ge7\lc\sqrt m\ge2$. Count: $\#Z=\sum_kn_k=N-\#K$ (Lemma~\ref{lem:cuts}(b)) and
$\#Z\le\frac{3m}2\cdot4L(d-1)\le9mt\lt$. $K\cap Z=\emptyset$ ($\dep\ge d+6$ versus $\dep\le d-1$). Discs: kept ones by Lemma~\ref{lem:keep}(a); filler ones lie in the
interiors of pieces with disjoint interiors, inside $O_-$; a kept and a filler point are at distance $\ge7>\frac12$.
\end{proof}

\begin{lemma}[Neumann identity]\label{lem:neumann}
$E'\in L^1(Q)$, $E'=0$ off $Q$, and $E'$ is a Neumann charge field on $Q$ with density $m$ and charges $Y$.
\end{lemma}
\begin{proof}
Supports: $\psi X$ in $\{\dep\ge d-1\}$, $V_k$ in $H_k$, $V^c_j$ in $H^c_j$, $\Psi_P$ in $P$, all in $Q$; each term is integrable. For $\varphi\in C^\infty_c(\C)$, by
Lemmas~\ref{lem:soft}(b), \ref{lem:flux}(a), \ref{lem:carry}, \ref{lem:piece},
\[
\int_QE'\cdot\nabla\varphi=2\pi\sum_K\varphi(p)+\int_Q(g-2\pi m)\varphi-\int\Bigl(\sum_kf_k\mathbf 1_{H_k}+\sum_jf^c_j\Bigr)\varphi+\sum_{P}2\pi\Bigl(\varphi(z_P)-\int_P\mu\varphi\Bigr).
\]
A.e.\ on $Q$: $\sum_kf_k\mathbf 1_{H_k}+2\pi\mu\mathbf 1_{\bigcup P_k}=\tilde g$ (on $T_k$: $f_k=\tilde g$; on $H'_k$: $2\pi(m-\mu_k)+2\pi\mu_k=2\pi m=\tilde g$; on $P_k\setminus H'_k$:
$2\pi\mu=\tilde g$; on $\Omega_{\rm c}$ all vanish), and $\sum_jf^c_j=g-\tilde g$. Hence the right side is $2\pi\sum_Y\varphi(y)-2\pi m\int_Q\varphi$.
\end{proof}

\begin{lemma}[kept singularities against new fields: an exact pairing with a sign]\label{lem:sign}
Let $Y_F=\sum_kV_k-\sum_{q\in D^\ast}W_q$ and $O_2=\{\dep>d+\frac12\}$. For every Lipschitz $\varphi$ with compact support in $O_2$,
\[
\int Y_F\cdot\nabla\varphi=-2\pi\sum_{q\in D^\ast}\avg_{\partial B(q,1/4)}\varphi .
\]
Consequently, for $p\in K$, $\int\nabla F_p\cdot Y_F=-2\pi\sum_q\avg_{\partial B(q,1/4)}F_p\ge0$, and $\int_{\Omega_+}E_\beta\cdot Y_F\le\int_{\Omega_+}\Phi\cdot Y_F$.
\end{lemma}
\begin{proof}
On $O_2\cap Q$: $\psi=1$, $\nabla\psi=0$, no carry cell is present, and $O_2\cap F\subset T$ (Lemma~\ref{lem:pinwheel}(c)); so
$\sum_kf_k\mathbf 1_{H_k}=\tilde g=g=2\pi\sum_{D^\ast}u_q$ there, and on $O_2\cap\Omega_{\rm c}$ both sides vanish. By Lemma~\ref{lem:T2}(b),
$\sum_k\int V_k\cdot\nabla\varphi=-2\pi\sum_q\int u_q\varphi$; by Lemma~\ref{lem:radial}(a), $\int W_q\cdot\nabla\varphi=2\pi(\avg_{\partial B(q,1/4)}\varphi-\int u_q\varphi)$. Subtract.
Here $F_p=f_{\alpha_p,1/4}(\cdot-p)$ denotes the continuous representative (Remark~\ref{rem:contrep}); it is the Lipschitz test function used, and the
a.e.\ gradient $\nabla F_p$ and all circle averages are those of the \texttt{truncLog} difference. $F_p$ is Lipschitz, $\le0$, with support $\bar B(p,\frac14)\subset\{\dep\ge d+\frac{23}4\}\subset O_2$, which gives the sign. Finally
$E_\beta=\Phi-\sum_K\nabla F_p$ (Lemma~\ref{lem:keep}(d)) and $\supp\nabla F_p\subset\Omega_+$.
\end{proof}

\begin{proposition}[energy of the screened field]\label{prop:energy}
Assume (K1)--(K4). With $E'_\rho=E'-\sum_{y\in Y}\nabla f_{\rho_y}(\cdot-y)$,
\[
G_Q(\rho)[E';Y]\le G_Q(\alpha)[E;C]+c_mn_\partial+\tfrac12S_d+32N_2+X_A+2V_{\rm flux}+\frac{12\pi^2}{\lc}+3\,\#Z\Bigl(\frac\pi2+e_T(m)\Bigr).
\]
\end{proposition}
\begin{proof}
Split $Q$ into $Q\cap\Omega_+$, $A$, $O_-$. Families $\{V_k\}$, $\{V^c_j\}$, $\{\Psi_P-\nabla f_{r_0}(\cdot-z_P)\}$ each have a.e.\ disjoint supports, so e.g.\
$|\sum_kV_k|^2=\sum_k|V_k|^2$ a.e.

\emph{On $Q\cap\Omega_+$}: $\psi=1$, carries and pieces vanish, filler discs are absent, so by Lemma~\ref{lem:soft}(a) $E'_\rho=E_\beta+Y_F$. With Lemma~\ref{lem:sign},
$\supp Y_F\cap\Omega_+\subset\{d<\dep\le d+8\}\subset R_d$, and $xy\le\frac12x^2+\frac12y^2$:
\[
\tfrac12\int_{\Omega_+}|E'_\rho|^2\le\tfrac12\int_{\Omega_+}|E_\beta|^2+\int_{\Omega_+}\Phi\cdot Y_F+\tfrac12\int_{\Omega_+}|Y_F|^2\le\tfrac12\int_{\Omega_+}|E_\beta|^2+\tfrac12S_d+\int_{\Omega_+}|Y_F|^2,
\]
and $\int_{\Omega_+}|Y_F|^2\le2\int_{\Omega_+}\sum_k|V_k|^2+2\int(4n)^2$.

\emph{On $A$}: no disc of $Y$ meets $A$, carries and pieces vanish: $E'_\rho=\psi X+\sum_kV_k$ and $\frac12\int_A|E'_\rho|^2\le X_A+\int_A\sum_k|V_k|^2$.

\emph{On $O_-$}: $\psi=0$ and $E'_\rho=\sum_kV_k+\sum_jV^c_j+\sum_P(\Psi_P-\nabla f_{r_0}(\cdot-z_P))$; by $\frac12|x+y+z|^2\le\frac32(|x|^2+|y|^2+|z|^2)$ and
Lemmas~\ref{lem:carry}, \ref{lem:piece}: $\frac12\int_{O_-}|E'_\rho|^2\le\frac32\int_{O_-}\sum_k|V_k|^2+\frac32\cdot3\cdot\frac{8\pi^2}{3\lc}+3\#Z(\frac\pi2+e_T)$.

Adding, the $V_k$ terms total at most $2V_{\rm flux}$, so
$\frac12\int_Q|E'_\rho|^2\le\frac12\int_Q|E_\beta|^2+\frac12S_d+32N_2+X_A+2V_{\rm flux}+\frac{12\pi^2}\lc+3\#Z(\frac\pi2+e_T)$.
Finally $\sum_Y\tau_m(\rho_y)=\sum_K\tau_m(\alpha_p)+\#Z\,\tau_m(r_0)\le\sum_K\tau_m(\alpha_p)$ ($r_0<1$), and Lemma~\ref{lem:keep}(c).
\end{proof}

\subsection{\texorpdfstring{Constants and the proof of Theorems~\ref{thm:I4int} and~\ref{thm:I4}}{Constants and the proofs of the screening theorems}}

\begin{definition}[constants]\label{def:constants}
\begin{align*}
\lambda_0&=\max\Bigl\{36,\ \frac8{m_-},\ \frac2{\sqrt{m_-}},\ 8\Bigl(9+\frac{13}{12}\cdot\frac{175}4\,\pi R_1^2+\frac1{36m_-}\Bigr)\Bigr\},\\
C_1&=\Bigl(\frac8{m_-}\Bigr)^2\Bigl(\frac{13}{12}b_1\sqrt{1400M_1}+\frac4\pi\Bigr)^2,
\end{align*}
and, for $m\in[m_-,m_+]$, with $\nu_t=\frac{33\pi}2a_2^2$, $\nu_S=4\pi b_2^2M_2$, $\gamma_t=4(144\pi^2m^2+1056\nu_t)$, $\gamma_S=4(1056\nu_S+64)$:
\[
K_S(m)=1008\gamma_S+80+64\nu_S,\qquad K_T(m)=1008\gamma_t+64\nu_t+48384\pi^2+\frac{\pi^2}3+27m\Bigl(\frac\pi2+e_T(m)\Bigr),
\]
$C_S=\max\{c_{m_+},\ \sup_mK_S(m),\ \sup_mK_T(m)\}$ (finite: continuous in $m$ on a compact interval).
\end{definition}

\begin{lemma}[the conditions (K1)--(K4) hold]\label{lem:Kholds}
If $\lc\ge\lambda_0$ and $C_1$ is as in Definition~\ref{def:constants}, then (K1)--(K4) hold for every $m\in[m_-,m_+]$.
\end{lemma}
\begin{proof}
(K2)--(K4) are in $\lambda_0$. (K1): divide by $2\pi w\lc$ and use $w\ge\lc\ge36$ (so $2+\frac3w\le\frac{25}{12}$, $\frac1{w\lc}\le\frac1{36\lc}$, $\frac{21\cdot25}{12}=\frac{175}4$,
$672\cdot\frac{25}{12}=1400$):
\[
\frac{\Theta(w)}{2\pi w\lc}\le\frac{9m+\frac{13}{12}\frac{175}4a_1+\frac1{36}}{\lc}+\frac{13}{12}b_1\sqrt{\frac{1400M_1}{C_1}}+\frac4{\pi\sqrt{C_1}}\le\frac m8+\frac m8,
\]
the first by $\lc\ge\lambda_0$ (note $a_1=m\pi R_1^2$), the second by the choice of $C_1$ ($m\ge m_-$). (A grid check over $\lc\in[\lambda_0,10^{10}\lambda_0]$,
$w\in[\lc,3\lc]$, $m\in[m_-,m_+]$ with $S_d$ maximal gives a worst ratio $\le0.971$.)
\end{proof}

\begin{proof}[Proof of Theorem~\ref{thm:I4int}]
By Lemma~\ref{lem:Kholds}, Lemmas~\ref{lem:points}, \ref{lem:neumann} and Proposition~\ref{prop:energy} apply. Insert Lemma~\ref{lem:ring}(b),(c),
Lemma~\ref{lem:flux}(b),(c), $S_d\le32\sigma$ and $\#Z\le9mt\lt$: $N_2\le\nu_tt+\nu_S\sigma$, $X_A\le64\sigma+32N_2$,
$\hat G\le\gamma_tt+\gamma_S\sigma+48\pi^2/\lc$, so the excess over $G_Q(\alpha)$ is at most
\[
c_mn_\partial+(80+64\nu_S)\sigma+1008\gamma_S\lc\sigma+1008\gamma_t\lc t+64\nu_tt+48384\pi^2+\frac{12\pi^2}{\lc}+27m\Bigl(\frac\pi2+e_T\Bigr)t\lt .
\]
Use $\sigma\le\lc\sigma$, $t\le t\lt$, $\lc t\le t\lt$, $1\le t\lt$, $12\pi^2/\lc\le\pi^2/3$: the excess is $\le c_mn_\partial+K_S\lc\sigma+K_Tt\lt\le C_S(\lc S/\lt+t\lt+n_\partial)$.
\end{proof}

\begin{remark}[numerical values]\label{rem:numbers}
$m_-=m_+=1$ (the jellium lower bound): $\lambda_0=3135.2$, $C_1=9.245\cdot10^5$, $\ell_0=25082$, $C_{\rm scr}=4.733\cdot10^8$, $C_S=4.27\cdot10^{10}$.
$[m_-,m_+]=[\frac1{3750},\frac1{150}]$ (the sphere note): $\lambda_0=30000$, $C_1=1.300\cdot10^{13}$, $\ell_0=2.4\cdot10^5$, $C_{\rm scr}=6.66\cdot10^{15}$, $C_S=1.11\cdot10^9$.
No attempt was made to optimize.
\end{remark}

\begin{remark}[three points of the construction, and the comparison with the Lectures]\label{rem:risks}
\emph{(1) Column masses after dropping charges near the cut.} The dropped charges are smeared into the discs $u_q$ (mass $\ge0$); the only signed part of $g$ is
$-\nabla\psi\cdot X$ on the unit ramp, whose integral over any $w\times9$ top strip is $\le\sqrt{2wS_d}+\frac\pi6n_D$, while a strip carries background mass
$\ge mw(d-1)\ge7mw\lc$. Lemma~\ref{lem:cuts}(d) quantifies it. \emph{(2) Cross terms with tiny kept radii.} They are computed exactly (Lemma~\ref{lem:sign}) and have
a favourable sign, because the new fields near kept points ($\sum V_k-\sum W_q$) have divergence $2\pi\sum\sigma_{q,1/4}\ge0$ there and $F_p\le0$: no lower bound on
$\alpha_p$ is needed. \emph{(3) Integer cuts versus cell averages.} The flux cells are subordinate to the columns: the average $\mu_k$ is taken after the cut and the
flux cell preserves the column mass, so the filler receives exactly the integer $n_k$. \emph{Choices made in the construction:} ring
$\{d-5\le\dep\le d+11\}$ (width $16$); kept zone $\dep\ge d+6$; frame $\dep\le d+8$; compensation height $\lc$; three carries of energy $O(1/\lc)$ each (oriented
across the common side); the $\lc$ used internally is $\frac18$ of the $\lc$ of the statement. Compared with Appendix~A of the Lectures: there is no Neumann problem, no trace, no
use of Lemma~A.1 or A.2 (Lemma~A.2 contains an error, see the errata record), no integer-mass partition left undone (it is carried out here), no corner analysis; the filler is explicit.
\end{remark}

%%%%%%%%%%%%%%%%%%%%%%%%%%%%%%%%%%%%%%%%%%%%%%%%%%%%%%%%%%%%%%%%%%%%%%%%%%%%%%%%%%%%%%%%%%%%%%%%%%%%%%%%%
\section{\texorpdfstring{Projection onto $T_N$}{Projection onto T\_N}}\label{sec:projection}

\textbf{Setting.} $N\ge1$, $\Gamma=\Gamma_N=\texttt{squareLattice}\,N=\sqrt N\,\Z[i]$ (covolume $N$), $T_N=\C/\Gamma$, and $Q=\sq(c,\sqrt N)$, a closed square whose
$\Gamma$-translates cover $\C$ and overlap in null sets (an a.e.\ fundamental domain; the translates of $\operatorname{int}Q$ are pairwise disjoint and
$(\gamma+\operatorname{int}Q)\cap Q=\emptyset$ for $\gamma\ne0$). The energy compared is \texttt{torusEnergy}$(\Gamma_N,a)$ of
Definition~\ref{def:objects}(g), the torus is $T_N$ itself (no reflection: the square $Q$ \emph{is} a fundamental domain of $\Gamma_N$; reflection is needed
only for the sphere note, \S\ref{sec:sphere}). Let $a:\mathrm{Fin}\,N\to\operatorname{int}Q$ be injective and $P=a+\Gamma$. The classes of the $a_i$ are pairwise distinct
(two points of $\operatorname{int}Q$ differ by less than $\sqrt N$ in each coordinate). $H=\texttt{torusPotential}\,\Gamma\,a$.

\begin{lemma}[periodization]\label{lem:period}
Let $E'$ be a Neumann charge field on $Q$ with density $1$ and charges $\{a_i\}$. Define $E^{\rm per}(x)=E'(x-\gamma)$ if $x-\gamma\in\operatorname{int}Q$
($\gamma\in\Gamma$ unique), $E^{\rm per}(x)=0$ otherwise. Then $E^{\rm per}$ is $\Gamma$-periodic, $E^{\rm per}=E'$ on $\operatorname{int}Q$, and
$E^{\rm per}\in\mathrm{elec}(P)$.
\end{lemma}
\begin{proof}
Periodicity and local integrability are clear (a compact set meets finitely many translates of $Q$), and $P$ is locally finite. For $\varphi\in C^\infty_c(\C)$ let
$\Phi=\sum_\gamma\varphi(\cdot+\gamma)$ (locally finite sum; smooth, periodic), $\chi\in C^\infty_c$ with $\chi=1$ near $Q$, and $\hat\varphi=\chi\Phi\in C_c^\infty$.
Then $\int E^{\rm per}\cdot\nabla\varphi=\sum_\gamma\int_{\operatorname{int}Q}E'(y)\cdot\nabla\varphi(y+\gamma)\dd y=\int_QE'\cdot\nabla\hat\varphi
=2\pi(\sum_i\hat\varphi(a_i)-\int_Q\hat\varphi)=2\pi(\sum_{p\in P}\varphi(p)-\int_\C\varphi)$, the last step by unfolding over $Q$.
\end{proof}

\begin{lemma}[periodic unfolding]\label{lem:unfold}
Let $Z\in L^1_{\rm loc}(\C;\C)$ be $\Gamma$-periodic with $\int Z\cdot\nabla\varphi=0$ for all $\varphi\in C^\infty_c(\C)$. Then $\int_QZ\cdot\nabla\psi=0$ for every
$\Gamma$-periodic Lipschitz $\psi$.
\end{lemma}
\begin{proof}
Let $\theta_0\in C_c^\infty$, $\theta_0\ge0$, $\theta_0>0$ on $Q$, and $\theta=\theta_0/\sum_\gamma\theta_0(\cdot+\gamma)$ (the denominator is smooth, periodic and positive), so
$\theta\in C^\infty_c$ and $\sum_\gamma\theta(\cdot+\gamma)\equiv1$. By Lemma~\ref{lem:moll} ($\nu=0$), $0=\int Z\cdot\nabla(\theta\psi)=\int\theta\,Z\cdot\nabla\psi+\int\psi\,Z\cdot\nabla\theta$.
Unfold both integrals over $Q$ (finite sums): by periodicity of $Z,\psi,\nabla\psi$ the first is $\int_QZ\cdot\nabla\psi\sum_\gamma\theta(\cdot+\gamma)=\int_QZ\cdot\nabla\psi$ and the second
is $\int_Q\psi Z\cdot\nabla\bigl(\sum_\gamma\theta(\cdot+\gamma)\bigr)=0$.
\end{proof}

\begin{lemma}[local form, $N$ charges]\label{lem:localform}
For each $i$ there are $s_i>0$ and $v_i\in C^\infty(B(a_i,s_i))$ with $H=v_i-\log|\cdot-a_i|$ on $B(a_i,s_i)\setminus\{a_i\}$. If $\bar B$ is a closed disc containing $a_i$ and
no other point of $P$, then $H+\log|\cdot-a_i|$ extends to a $C^\infty$ function on a neighbourhood of $\bar B$.
\end{lemma}
\begin{proof}
$H$ satisfies \texttt{IsTorusPotential} (\texttt{isTorusPotential\_torusPotential}, distinct classes, covolume $N$). Take $s_i<\dist(a_i,P\setminus\{a_i\})$. On the punctured disc
$h=H+\log|\cdot-a_i|-\frac\pi2|\cdot-a_i|^2$ is smooth and harmonic ($\Delta H=2\pi$, $\Delta\log=0$) and has a finite limit at $a_i$, so it is bounded near $a_i$ and
extends harmonically (Input~\ref{in:removable}); put $v_i=h+\frac\pi2|\cdot-a_i|^2$. The second claim: off $a_i$ the function is smooth near $\bar B$ (no other charge), and
near $a_i$ it equals $v_i$. (Lean: the one-charge case is \texttt{Torus.exists\_local\_form\_ball}; the proof is local.)
\end{proof}

\begin{lemma}[truncated periodic potential]\label{lem:Htilde}
Let $\rho_i>0$ with $\bar B(a_i,\rho_i)\subset\operatorname{int}Q$ and $\rho_i+\rho_j\le|a_i-a_j|$ for $i\ne j$, and $\tilde H=H-\sum_{i,\gamma}f_{\rho_i}(\cdot-a_i-\gamma)$. Then $\tilde H$ is $\Gamma$-periodic and
Lipschitz, and a.e.\ on $Q$: $\nabla\tilde H=\nabla H-\sum_i\nabla f_{\rho_i}(\cdot-a_i)=:(\nabla H)_\rho$; in particular $(\nabla H)_\rho\in L^\infty(Q)$.
\end{lemma}
\begin{proof}
On $\bar B(a_i+\gamma,\rho_i)$, $\tilde H=(H+\log|\cdot-a_i-\gamma|)-\log\rho_i$, smooth on a neighbourhood (Lemma~\ref{lem:localform}, translated; the disc contains no other
charge). On the closed set $\C\setminus\bigcup B(a_i+\gamma,\rho_i)$, $\tilde H=H$ is smooth on a neighbourhood. $\tilde H$ is continuous ($f_\rho=0$ on the circles).
$|\nabla\tilde H|\le\Lambda$ on both pieces (periodicity and compactness). A segment meets each circle in at most two points, so it splits into finitely many closed
subsegments, each in one piece; integrating $\nabla\tilde H$ along them gives $|\tilde H(x)-\tilde H(y)|\le\Lambda|x-y|$. On $Q$ only $\gamma=0$ contributes.
\end{proof}

\begin{theorem}[\texttt{torusEnergy\_le\_locEnergy}]\label{thm:I5}
Let $a$ be as above, $\rho_i>0$ with $\bar B(a_i,\rho_i)\subset\operatorname{int}Q$ and $\rho_i+\rho_j\le|a_i-a_j|$ ($i\ne j$), and $E'$ a Neumann charge field on $Q$
with density $1$ and charges $\{a_i\}$. Then
\[
\texttt{torusEnergy}(\Gamma_N,a)\ \le\ G_Q(\rho)[E';\{a_i\}]\qquad(m=1).
\]
\end{theorem}
\begin{proof}
If $E'_\rho\notin L^2(Q)$ the right side is $+\infty$. Otherwise:

\emph{Step 1.} $E^{\rm per}\in\mathrm{elec}(P)$ (Lemma~\ref{lem:period}) and $\nabla H\in\mathrm{elec}(P)$ (Input~\ref{in:main}); hence $Z:=E^{\rm per}-\nabla H$ is periodic,
$L^1_{\rm loc}$, and $\int Z\cdot\nabla\varphi=0$ for $\varphi\in C^\infty_c$.

\emph{Step 2.} On $\operatorname{int}Q$, $E'_\rho=Z+(\nabla H)_\rho=Z+\nabla\tilde H$ a.e., so $Z\in L^2(Q)$, and Lemma~\ref{lem:unfold} with $\psi=\tilde H$
(Lemma~\ref{lem:Htilde}) gives $\int_QZ\cdot\nabla\tilde H=0$. Therefore $\frac12\int_Q|E'_\rho|^2=\frac12\int_Q|(\nabla H)_\rho|^2+\frac12\int_Q|Z|^2$ and
$G_Q(\rho)[E']\ge G_Q(\rho)[\nabla H]$ (same $\tau$-terms).

\emph{Step 3: $G_Q(\rho)[\nabla H]=G_Q(\eta)[\nabla H]$ for small $\eta$.} Let $d_i=\dist(a_i,\C\setminus\operatorname{int}Q)-\rho_i>0$, $s_i$ as in Lemma~\ref{lem:localform},
$\eta_0=\min_i\min\{\rho_i,d_i,s_i\}$ and $0<\eta<\eta_0$. Apply Proposition~\ref{prop:mono} and Corollary~\ref{cor:mono}(b) to $\nabla H\in\mathrm{elec}(P)$ ($U=\C$), $D=Q$,
$\alpha\equiv\eta$ on $P$, $\beta=\rho_i$ at $a_i$ and $\eta$ elsewhere: the changed points are the $a_i$, $\bar B(a_i,\rho_i)\subset Q$, $(\nabla H)_\eta\in L^2(Q)$ (bounded, by
Lemma~\ref{lem:localform}), and the disjointness holds: $\rho_i+\rho_j\le|a_i-a_j|$, and for $q=a_j+\gamma$, $\gamma\ne0$, $q\notin\operatorname{int}Q$, so
$|a_i-q|\ge\rho_i+d_i>\rho_i+\eta$. So $G_Q(\beta)=G_Q(\eta)$. On $Q$ the truncations at $a_j+\gamma$ ($\gamma\ne0$) vanish (those discs miss $Q$) and $P\cap Q=\{a_i\}$,
so $G_Q(\beta)[\nabla H;P]=G_Q(\rho)[\nabla H;\{a_i\}]$, and, as $\nabla H-\nabla f_\eta(\cdot-a_i)=\nabla v_i$ on $B(a_i,\eta)$,
\[
G_Q(\eta)[\nabla H;P]=\tfrac12\int_{Q\setminus\bigcup_iB(a_i,\eta)}|\nabla H|^2+\tfrac12\sum_i\int_{B(a_i,\eta)}|\nabla v_i|^2+N(\pi\log\eta-\pi^2\eta^2).
\]

\emph{Step 4: the limit.} $|\nabla H|^2\mathbf 1_{\C\setminus\bigcup_{p\in P}B(p,\eta)}$ is periodic and bounded, so its integrals over the a.e.\ fundamental domains $T$ and $Q$
agree (\texttt{torusDomain\_spec}), and $Q\setminus\bigcup_PB(p,\eta)=Q\setminus\bigcup_iB(a_i,\eta)$. Hence the expression $J(\eta)$ in the definition of \texttt{torusEnergy}
satisfies $J(\eta)=G_Q(\eta)-\frac12\sum_i\int_{B(a_i,\eta)}|\nabla v_i|^2+N\pi^2\eta^2\to G_Q(\rho)[\nabla H]$ as $\eta\downarrow0$. The limit along $\mathcal N[>]0$ is
unique, so $\texttt{torusEnergy}(\Gamma_N,a)=G_Q(\rho)[\nabla H]\le G_Q(\rho)[E']$.
\end{proof}

\begin{remark}
Equality holds iff $Z=0$, i.e.\ $E'=\nabla H$ on $Q$. Step 4 is the $N$-point generalization of \texttt{TriangularEnergyLimit.tendsto\_puncture} (one charge);
with Step 3 it shows $\texttt{torusEnergy}=G_Q(\rho)[\nabla H]$ for \emph{every} admissible $\rho$, which is the truncation identity of the sphere note (Lemma~\texttt{lem:truncation}) in this
setting.
\end{remark}

%%%%%%%%%%%%%%%%%%%%%%%%%%%%%%%%%%%%%%%%%%%%%%%%%%%%%%%%%%%%%%%%%%%%%%%%%%%%%%%%%%%%%%%%%%%%%%%%%%%%%%%%%
\section{Assembly: the jellium lower bound and the configuration lower bound}\label{sec:assembly}

\begin{lemma}[finite neutral energy]\label{lem:finiteW}
Let $E\in\mathrm{elec}(C)$ with $\Wneu(E)<+\infty$. Then (a) $n(\ell)/\ell^2\to1$; (b) there are $M\ge0$, $\ell_1$ with $F^{\square_\ell}(E)\le M\ell^2$ for all real $\ell\ge\ell_1$;
(c) with $A=4\pi M+4c_1$ there is $\ell_2$ with $\int_{\square_\ell}|E_{1/4}|^2\le A\ell^2$ for $\ell\ge\ell_2$; (d) $n_\partial(\ell):=\#(C\cap\square_\ell\setminus\square_{\ell-4\ell^{3/4}})=o(\ell^2)$.
\end{lemma}
\begin{proof}
(a) Definition of \texttt{neutralEnergy}. (b) An EReal $\limsup$ below $+\infty$ is below some real $M$, and then the sequence is eventually $\le M$.
(c) Corollary~\ref{cor:EB} and $n(\ell)\le2\ell^2$ eventually. (d) $n_\partial(\ell)=n(\ell)-n(\ell-4\ell^{3/4})$ and $(\ell-4\ell^{3/4})/\ell\to1$.
\end{proof}

\begin{proposition}[the cube energy from the torus bound]\label{prop:I6}
Assume the statement of Input~\ref{in:R1}. Let $E\in\mathrm{elec}(C)$ with $\Wneu(E)\ne+\infty$ and $\varepsilon>0$. Then for all sufficiently large $N\in\N$,
\[
N\Bigl(\frac{W_\triangle}{2\pi}-\varepsilon\Bigr)\le F^{\square_{\sqrt N}}(E).
\]
\end{proposition}
\begin{proof}
Let $\ell_0,C_{\rm scr},C_S$ be the constants of Theorem~\ref{thm:I4} for $m_-=m_+=1$, and $M,A$ as in Lemma~\ref{lem:finiteW}. For $N\in\N$ put $\ell=\sqrt N$ and apply
Theorem~\ref{thm:I4} with $c=0$, $t=\ell$, $m=1$, $\lc=\ell^{1/2}$, $\lt=\ell^{3/4}$, the given $C$, $E$, and $\alpha=\tilde r=\texttt{cubeRadius}\,C\,\ell$. For $N$ large the
hypotheses hold: (H1) $\ell_0\le\ell^{1/2}\le\ell^{3/4}$, $8\ell^{3/4}\le\ell$ ($N\ge\max\{\ell_0^4,8^8\}$), $N=1\cdot t^2$; (H2) Lemma~\ref{lem:norm}(a);
(H3) $0<\tilde r\le\frac14$, and if $\dep(p)<1$ then either $p\notin\square_\ell$ or $\dist(p,\partial\square_\ell)=\dep(p)<1$ by \eqref{eq:depth-sets}, so $\tilde r_p=\frac14$;
(H4) on $\{\dep\ge\ell^{3/4}\}$ ($\ge2$), $\tilde r_p=r_p\le\frac14|p-q|$, so $\tilde r_p+\tilde r_q\le|p-q|$; (H5) $F^{\square_\ell}<\infty$ (Lemma~\ref{lem:finiteW}(b));
(H6) $S\le\int_{\square_\ell}|E_{1/4}|^2\le A\ell^2\le\ell^{9/4}/C_{\rm scr}$ once $\ell^{1/4}\ge AC_{\rm scr}$. We get $Y,\rho,E'$; enumerate $Y$ as an injective
$a:\mathrm{Fin}\,N\to\operatorname{int}\square_\ell$. By Theorem~\ref{thm:I5} ($\sq(0,\sqrt N)=\square_\ell$), Input~\ref{in:R1} ($N\ge2$, distinct classes) and
Lemma~\ref{lem:norm}(b),
\[
NW_\triangle\le\texttt{torusEnergy}(\Gamma_N,a)\le G_{\square_\ell}(\rho)[E']\le2\pi F^{\square_\ell}(E)+C_S\bigl(A\ell^{7/4}+\ell^{7/4}+n_\partial(\ell)\bigr),
\]
using $\lc S/\lt\le\ell^{1/2}A\ell^2/\ell^{3/4}$ and $t\lt=\ell^{7/4}$. Divide by $2\pi N=2\pi\ell^2$: $F^{\square_\ell}/N\ge W_\triangle/2\pi-\frac{C_S}{2\pi}\bigl((A+1)\ell^{-1/4}+n_\partial(\ell)/\ell^2\bigr)$,
and the bracket tends to $0$ (Lemma~\ref{lem:finiteW}(d)). 
\end{proof}

\begin{theorem}[the jellium lower bound; \texttt{div\_two\_pi\_triangularEnergy\_le\_neutralEnergy}]\label{thm:R2}
For every $C$ and $E\in\mathrm{elec}(C)$: $W_\triangle/2\pi\le\Wneu(E)$.
\end{theorem}
\begin{proof}
Trivial if $\Wneu(E)=+\infty$. Otherwise $n(\ell)/\ell^2\to1$ and $\Wneu(E)=\limsup_{\ell\to\infty}\ell^{-2}F^{\square_\ell}(E)$. Given $\varepsilon>0$, by
Proposition~\ref{prop:I6} the inequality $\ell^{-2}F^{\square_\ell}\ge W_\triangle/2\pi-\varepsilon$ holds at $\ell=\sqrt N$ for all large $N$, hence frequently as $\ell\to\infty$
(every $\ell_*$ is exceeded by some $\sqrt N$); \texttt{le\_limsup\_of\_frequently\_le} gives $\Wneu(E)\ge W_\triangle/2\pi-\varepsilon$. Let $\varepsilon\downarrow0$.
\end{proof}

\begin{lemma}[$E_\triangle$ is a distributional gradient]\label{lem:Etrigrad}
$H_\triangle=\texttt{torusPotential}\,\Lambda_\triangle\,0\in L^1_{\rm loc}$ and $\int H_\triangle\nabla\varphi=-\int\varphi E_\triangle$ for every real $\varphi\in C^\infty_c$.
\end{lemma}
\begin{proof}
By the local form (\texttt{Torus.exists\_local\_form\_ball}, radius $1$, as nonzero lattice vectors have norm $\ge1$) and periodicity, near each $p\in\Lambda_\triangle$,
$|H_\triangle|\le C_0+\log(1/|x-p|)$ and $|\nabla H_\triangle|\le C_0+1/|x-p|$; both are $L^1_{\rm loc}$, and $H_\triangle$ is smooth off $\Lambda_\triangle$. Fix $\varphi$ and the finitely many
$p_1,\dots,p_k\in\Lambda_\triangle$ within distance $1$ of $\supp\varphi$. For $0<\epsilon<\frac12$, the divergence theorem for the compactly supported field $H_\triangle\varphi\,e$
($e\in\{1,i\}$) on $\C\setminus\bigcup_j\bar B(p_j,\epsilon)$ (\texttt{PolarDivergence.integral\_exterior\_div}) gives
$\int_{\C\setminus\bigcup B(p_j,\epsilon)}(H_\triangle\nabla\varphi+\varphi\nabla H_\triangle)=O\bigl(\sum_j\epsilon(C_0+\log\frac1\epsilon)\|\varphi\|_\infty\bigr)\to0$. Dominated convergence yields
$\int H_\triangle\nabla\varphi=-\int\varphi\nabla H_\triangle$, and $E_\triangle=\nabla H_\triangle$ off the null set $\Lambda_\triangle$.
\end{proof}

\begin{theorem}[the configuration lower bound; \texttt{configEnergy\_triangularLattice}]\label{thm:R3}
$\mathbb W(\Lambda_\triangle)=W_\triangle/2\pi$ and $\mathbb W(\Lambda_\triangle)\le\mathbb W(C)$ for every $C$.
\end{theorem}
\begin{proof}
For every $C$, each term of the infimum defining $\mathbb W(C)$ is $\ge W_\triangle/2\pi$ by Theorem~\ref{thm:R2} (and $\inf\emptyset=+\infty$), so $\mathbb W(C)\ge W_\triangle/2\pi$. For
$\Lambda_\triangle$, $E_\triangle$ is admissible (Input~\ref{in:Etri} and Lemma~\ref{lem:Etrigrad}), so $\mathbb W(\Lambda_\triangle)\le\Wneu(E_\triangle)=W_\triangle/2\pi$. Hence equality, and then
$\mathbb W(\Lambda_\triangle)=W_\triangle/2\pi\le\mathbb W(C)$.
\end{proof}

%%%%%%%%%%%%%%%%%%%%%%%%%%%%%%%%%%%%%%%%%%%%%%%%%%%%%%%%%%%%%%%%%%%%%%%%%%%%%%%%%%%%%%%%%%%%%%%%%%%%%%%%%
\section{The screening input of the sphere note from the screening theorem}\label{sec:sphere}

References are to the sphere note: Definition~6.1 (\texttt{def:neumann}), Lemma~6.2 (\texttt{lem:reflection}),
Input~6.3 (\texttt{in:SCR}) and Lemma~\texttt{lem:count}.

\begin{lemma}[Neumann projection with reflection]\label{lem:reflproj}
Let $Q$ be a closed square of side $t$, $\bar m>0$ with $\bar n=\bar mt^2\in\N$, $Y\subset\operatorname{int}Q$ with $\#Y=\bar n$, radii $\rho$ with $\bar B(y,\rho_y)\subset\operatorname{int}Q$ and $\rho_y+\rho_{y'}\le|y-y'|$ for $y\ne y'$, and $E'$ a Neumann charge field on $Q$ with density $\bar m$ and charges $Y$. Then $W^{\rm Neu}_Q(Y;\bar m)\le G_Q(\rho)[E';Y]$ ($\tau_{\bar m}$ in $G$).
\end{lemma}
\begin{proof}
\emph{$\bar m=1$.} Let $\hat Q$ be the square of side $2t$ formed by $Q$ and its reflections $\mathcal R_1Q,\mathcal R_2Q,\mathcal R_1\mathcal R_2Q$ across two adjacent sides. On
$\mathcal R Q$ put $\hat E=\mathcal R_{\rm lin}E'\circ\mathcal R$ (linear part of the isometry $\mathcal R$). For $\varphi\in C^\infty_c$, the change of variables $y=\mathcal Rx$ gives
$\int_{\mathcal RQ}\hat E\cdot\nabla\varphi=\int_QE'\cdot\nabla(\varphi\circ\mathcal R)=2\pi(\sum_Y\varphi(\mathcal Ry)-\int_Q\varphi\circ\mathcal R)$. Summing over the four copies, $\hat E$ is a
Neumann charge field on $\hat Q$ with density $1$ and the $4\bar n$ charges $\hat Y=\bigcup\mathcal R(Y)$ (distinct, in $\operatorname{int}\hat Q$, with radii satisfying $\hat\rho_y+\hat\rho_{y'}\le|y-y'|$: within one copy by hypothesis, across copies because the discs lie in the interiors of different copies), and
$G_{\hat Q}(\hat\rho)[\hat E]=4G_Q(\rho)[E']$. Since $2t=\sqrt{4\bar n}$, Theorem~\ref{thm:I5} gives $\texttt{torusEnergy}(\Gamma_{4\bar n},\hat Y)\le4G_Q(\rho)[E']$, and
Lemma~6.2(iii) of the sphere note (its $\tilde Y$ is $\hat Y$ modulo $2t\Z[i]$; \texttt{torusEnergy} depends only on the set of classes) gives $4W^{\rm Neu}_Q(Y;1)=\texttt{torusEnergy}$.

\emph{General $\bar m$.} With $\varsigma=\sqrt{\bar m}$, $\hat E(x)=\varsigma^{-1}E'(x/\varsigma)$ is a Neumann charge field on $\varsigma Q$ with density $1$ and charges $\varsigma Y$, and
(using $f_{\varsigma r}(\varsigma y)=f_r(y)$, $\tau_1(\varsigma r)=\tau_{\bar m}(r)+\frac\pi2\log\bar m$) $G_{\varsigma Q}(\varsigma\rho)[\hat E]=G_Q(\rho)[E']+\frac\pi2\bar n\log\bar m$. By the case
$\bar m=1$ and Lemma~6.2(iv), $W^{\rm Neu}_Q(Y;\bar m)=W^{\rm Neu}_{\varsigma Q}(\varsigma Y;1)-\frac\pi2\bar n\log\bar m\le G_Q(\rho)[E']$.
\end{proof}

\begin{proposition}[Input 6.3 of the sphere note from Theorem~\ref{thm:I4}]\label{prop:sphere}
Let $0<m_-\le m_+\le\frac1{150}$, and $\ell_0,C_{\rm scr},C_S$ the constants of Theorem~\ref{thm:I4} for $(m_-,m_+)$. Then Input~6.3 of the sphere note holds with these $\ell_0$,
$C_{\rm scr}$ and with $C_S$ replaced by $(1+\frac{48}\pi)C_S$ (for every $\ell\ge\ell_0$; the extra condition $\ell\ge4m_+/m_-$ is not needed).
\end{proposition}
\begin{proof}
Given the data of Input~6.3 ($Q$, $\bar m$, $\ell\le\tilde\ell\le t/8$, $X$ $1$-separated, $w$ with $w^\natural=w-\sum f_{1/4}(\cdot-x_i)\in H^1(Q)$ and
$-\Delta w=2\pi(\sum_{x_i\in Q^{1/2}}\delta_{x_i}-\bar m)$ in $\mathcal D'((Q^{1/2})^\circ)$), apply Theorem~\ref{thm:I4} with $C=X$, $E=\nabla w$, $\alpha\equiv\frac14$,
$\lc=\ell$, $\lt=\tilde\ell$, $m=\bar m$: $U=(Q^{1/2})^\circ$ and test functions supported in $U$ do not see charges outside $U$, so (H2) holds; (H3) trivially;
(H4) $\frac14+\frac14\le1\le|x_i-x_j|$; (H5) $E_\alpha=\nabla w^\natural\in L^2(Q)$; (H6) is Input~6.3's screenability ($Q^{\tilde\ell/2}\setminus Q^{2\tilde\ell}=\{\tilde\ell/2\le\dep<2\tilde\ell\}$ by
\eqref{eq:depth-sets}). We get $Y,\rho,E'$, and by Lemma~\ref{lem:reflproj},
$W^{\rm Neu}_Q(Y;\bar m)\le G_Q(\frac14)[\nabla w;X]+C_S(\ell S/\tilde\ell+t\tilde\ell+n_\partial)$. Now
$G_Q(\frac14)[\nabla w;X]=\frac12\int_Q|\nabla w^\natural|^2+\pi p\log\frac14-\frac{\pi^2\bar m}{16}p\le\frac12\int_Q|\nabla w^\natural|^2+\pi p\log\frac14-2\pi\bar m\sum_i\int_Qf_{1/4}(x-x_i)\dd x$
since $\int_Qf_{1/4}(\cdot-x_i)\le\pi/32$. Finally Lemma~\texttt{lem:count} of the sphere note gives $n_\partial\le\frac4\pi|\{-\frac12<\dep<2\tilde\ell+\frac12\}|\le\frac4\pi(4\tilde\ell+2)2t\le\frac{48}\pi t\tilde\ell$.
\end{proof}

\begin{remark}[consequence for the sphere note]
Proposition~\ref{prop:sphere} proves Input~6.3 of the sphere note outright, for $\ell\ge\ell_0$; the filler Sublemmas~6.16--6.17 (\texttt{lem:filler}, \texttt{lem:filler-energy}) of the sphere note are
then not needed for Input~6.3.
\end{remark}

%%%%%%%%%%%%%%%%%%%%%%%%%%%%%%%%%%%%%%%%%%%%%%%%%%%%%%%%%%%%%%%%%%%%%%%%%%%%%%%%%%%%%%%%%%%%%%%%%%%%%%%%%
\section{Summary: dependencies and Lean statements}\label{sec:summary}

\textbf{Dependency chain.}
Bridge to Definition~11.4: Lemma~\ref{lem:moll} $\to$ Prop.~\ref{prop:pairing} $\to$ Prop.~\ref{prop:mono} $\to$ Cor.~\ref{cor:EB}; Lemma~\ref{lem:localcount}, Lemma~\ref{lem:grid} $\to$
Prop.~\ref{prop:density} $\to$ Thm.~\ref{thm:R2prime}.
Jellium lower bound: Thm.~\ref{thm:I4} (via \S4) $+$ Thm.~\ref{thm:I5} $+$ Theorem~1.2 of the preprint (Input~\ref{in:R1}) $\to$ Prop.~\ref{prop:I6} $\to$ Thm.~\ref{thm:R2}.
Configuration lower bound: Thm.~\ref{thm:R2} $+$ \texttt{triangularField\_spec} $+$ Lemma~\ref{lem:Etrigrad} $\to$ Thm.~\ref{thm:R3}.

\textbf{Lean statements.} Lemma~\ref{lem:norm}(b) is \texttt{two\_pi\_mul\_cubeEnergy}; Cor.~\ref{cor:mono} is
\texttt{truncInteraction\_nonpos} and \texttt{truncInteraction\_eq\_zero}; Cor.~\ref{cor:EB} is \texttt{half\_lintegral\_quarter\_le};
Prop.~\ref{prop:density} is \texttt{tendsto\_ncard\_cube\_div\_sq}; Thm.~\ref{thm:R2prime} is \texttt{jelliumEnergy\_eq\_neutralEnergy};
Thm.~\ref{thm:I4} is \texttt{exists\_screened\_neumannChargeField}; Thm.~\ref{thm:I5} is \texttt{torusEnergy\_le\_locEnergy};
Thm.~\ref{thm:R2} is \texttt{div\_two\_pi\_triangularEnergy\_le\_neutralEnergy}; Thm.~\ref{thm:R3} is \texttt{configEnergy\_triangularLattice}.
The pairing, the rectangle solver, the radial fields and the averaging lemma (Prop.~\ref{prop:pairing}, Lemmas~\ref{lem:T2}, \ref{lem:radial}, \ref{lem:avg}, \ref{lem:moll}) are
auxiliary results of the formalization; the columns, filler and assembly (Lemmas~\ref{lem:pinwheel}--\ref{lem:piece},
\ref{lem:goodcut}--\ref{lem:ring}, \ref{lem:points}--\ref{lem:Kholds}, Prop.~\ref{prop:energy}) prove Theorem~\ref{thm:I4}, \S\ref{sec:projection} proves Theorem~\ref{thm:I5},
and \S\ref{sec:assembly} gives the assembly.

\textbf{Completeness.} The proofs in \S\S1--7 are complete. The inputs are: Theorem~1.2 of the preprint (Input~\ref{in:R1}), \texttt{triangularField\_spec}, the facts of
Input~\ref{in:main}, Mathlib, Input~\ref{in:removable} (classical), and, for \S7 only, Lemma~6.2 of the sphere note.

\textbf{Remarks on the Lean formalization (not mathematical gaps).} (i) \texttt{gradient} of a Lipschitz function is the a.e.\ gradient (Rademacher), used in Lemma~\ref{lem:moll}.
(ii) The $N$-charge local form (Lemma~\ref{lem:localform}) generalizes \texttt{Torus.exists\_local\_form\_ball}. (iii) The closed square $Q$ is an \emph{a.e.} fundamental domain
(\texttt{IsAddFundamentalDomain} allows null overlaps). (iv) Energies are lower Lebesgue integrals cast to EReal; one works with real energies under the finiteness hypotheses and
converts once. (v) The filler and column cuts use \texttt{intermediate\_value\_Icc} on continuous cumulative masses; strict monotonicity
(density $\ge m/2$) gives the inverse in Lemma~\ref{lem:fillpart}.

\textbf{Errata.} The route avoids Lemma~11.2, Prop.~11.3, Lemma~11.9 and Appendix~A of the Lectures, whose printed statements contain defects
(see the errata record of the repository).

\begin{thebibliography}{9}
\bibitem{Preprint} OpenAI, \emph{Triangular minimality for planar Coulomb renormalized energy}, preprint, 23~September 2026.
\bibitem{Lectures} S.~Serfaty, \emph{Lectures on Coulomb and Riesz gases}, arXiv:2407.21194, 2024.
\end{thebibliography}

\end{document}
